% !TEX program = pdflatex
% !TEX encoding = UTF-8 Unicode
\documentclass[12pt,a4paper]{article}
\usepackage[utf8]{vietnam}
\usepackage{amsmath,amssymb}
\usepackage{tabularx}
\usepackage{tikz}
\usepackage{xcolor}
\usepackage[top=0.45cm, bottom=0.45cm, left=0.6cm, right=0.6cm]{geometry}
\usepackage[dethi]{ex_test}

\setlength{\parskip}{0pt}
\setlength{\parindent}{0pt}

% Dòng kẻ chấm mờ đều đặn chuẩn viết tay (khoảng cách 19.5pt)
\newcommand{\dongke}[1]{%
  \begingroup%
  \setlength{\parskip}{0pt}%
  \baselineskip=19.5pt%
  \vspace{1mm}%
  \foreach \i in {1,...,#1}{%
    \noindent{\color{gray!60}\dotfill}\par%
  }%
  \vspace{0.8mm}%
  \endgroup%
}

% Định nghĩa ô bài tập tự luận trong bảng
\newcommand{\oBaiTap}[3][5]{%
  \vspace{0.8mm}%
  \noindent\textbf{#2} #3\par%
  \vspace{0.8mm}%
  \noindent\textit{Bài giải:}%
  \dongke{#1}%
}

% Câu trắc nghiệm 4 đáp án dàn hàng ngang
\newcommand{\cauTN}[6]{%
  \noindent\textbf{Câu #1.} #2\par\vspace{1mm}%
  \noindent%
  \makebox[0.25\linewidth][l]{\textbf{A.} #3}%
  \makebox[0.25\linewidth][l]{\textbf{B.} #4}%
  \makebox[0.25\linewidth][l]{\textbf{C.} #5}%
  \makebox[0.25\linewidth][l]{\textbf{D.} #6}\par\vspace{1.5mm}%
}

\begin{document}
\pagestyle{empty}


% #############################################################################
% ========================= TOÁN 8 - TUẦN 03 - BUỔI 02: MÃ ĐỀ 101 ============
% ========================= MẶT TRƯỚC (TRANG 1) ===============================
% #############################################################################

\noindent
\begin{tabularx}{\linewidth}{|X|c|}
\hline
\textbf{PHIẾU KIỂM TRA ĐỊNH KỲ TOÁN 8 - TUẦN 03 - BUỔI 02} & \textbf{Mã đề: 101} \\
\textit{Chủ đề: Ôn tập Chương 1 - Đa thức} & \textit{Thời gian: 45 phút} \\
\hline
\end{tabularx}

\vspace{1mm}

\noindent
\begin{tabularx}{\linewidth}{|X|}
\hline
\rule[-2mm]{0pt}{7.5mm}\textbf{Họ và tên:} \textbf{\large Đào Viết Tuấn Anh} \hfill \textbf{Mã HS:} \texttt{TUANANH2411} \quad \\
\hline
\end{tabularx}

\vspace{1mm}

\noindent
\begin{tabularx}{\linewidth}{|p{3.2cm}|X|}
\hline
\centering\textbf{ĐIỂM SỐ} & \centering\textbf{LỜI NHẬN XÉT CỦA THẦY CÔ} \tabularnewline
\hline
\rule{0pt}{1.3cm} & \tabularnewline
\hline
\end{tabularx}

\vspace{1.2mm}

\noindent\textbf{BẢNG ĐIỀN ĐÁP ÁN TRẮC NGHIỆM (Mỗi câu đúng 0,5 điểm):}
\vspace{0.8mm}

\noindent
\renewcommand{\arraystretch}{1.35}
\begin{tabularx}{\linewidth}{|c|*{5}{>{\centering\arraybackslash}X|}}
\hline
\textbf{Câu} & \textbf{1} & \textbf{2} & \textbf{3} & \textbf{4} & \textbf{5} \\
\hline
\textbf{Đáp án} & & & & & \\
\hline
\end{tabularx}
\renewcommand{\arraystretch}{1.0}

\vspace{1.5mm}

\noindent\textbf{\large PHẦN I. TRẮC NGHIỆM KHÁCH QUAN (2,5 điểm)}
\vspace{1mm}

\cauTN{1}{Bậc của đa thức $M = 5x^3 y^2 - 2x^2 y^3 + 3x^4 y - 5x^3 y^2 + 1$ là:}{$4$}{$5$}{$6$}{$3$}

\cauTN{2}{Đơn thức nào sau đây đồng dạng với đơn thức $-4x^3 y^2$?}{$2x^2 y^3$}{$-4x^3 y$}{$7x^3 y^2$}{$5x^2 y^2$}

\cauTN{3}{Kết quả của phép nhân đa thức $(2x - 3)(x + 2)$ là:}{$2x^2 + x - 6$}{$2x^2 - x - 6$}{$2x^2 + 7x - 6$}{$2x^2 - 6$}

\cauTN{4}{Kết quả của phép chia $(18x^4 y^3 - 12x^3 y^2) : 6x^2 y^2$ là:}{$3x^2 y - 2xy$}{$3x^2 y - 2$}{$3x^2 y^2 - 2x$}{$3x^2 y - 2x$}

\cauTN{5}{Giá trị của biểu thức $P = x^2 - 6xy + 9y^2$ tại $x = 16$ và $y = 2$ là:}{$100$}{$144$}{$64$}{$400$}

\vspace{1mm}

\noindent\textbf{\large PHẦN II. TỰ LUẬN (7,5 điểm)}

\noindent\textbf{Câu 1 (2,5 điểm). Thực hiện phép tính với đa thức:}
\vspace{0.8mm}

\noindent
\begin{tabularx}{\linewidth}{|X|X|}
\hline
\multicolumn{2}{|p{0.97\linewidth}|}{Cho hai đa thức: $A = 3x^2 y - 2xy^2 + 5xy - 1$ \ và \ $B = 2x^2 y + 2xy^2 - 3xy + 4$.} \\
\hline
\oBaiTap[7]{a)}{Tính $C = A + B$:} & 
\oBaiTap[7]{b)}{Tính $D = A - B$:} \\
\hline
\oBaiTap[5]{c)}{Tìm bậc của đa thức $C$:} & 
\oBaiTap[5]{d)}{Tính giá trị của đa thức $C$ tại $x = 1; y = -1$:} \\
\hline
\end{tabularx}


\newpage

% #############################################################################
% ========================= TOÁN 8 - TUẦN 03 - BUỔI 02: MÃ ĐỀ 101 ============
% ========================= MẶT SAU (TRANG 2) =================================
% #############################################################################

\noindent
\begin{tabularx}{\linewidth}{|X|c|}
\hline
\textbf{PHẦN II. TỰ LUẬN (Tiếp theo) - TOÁN 8 TUẦN 03 BUỔI 02} & \textbf{Mã đề: 101} \\
\hline
\end{tabularx}

\vspace{1mm}

\noindent
\begin{tabularx}{\linewidth}{|X|}
\hline
\rule[-2mm]{0pt}{7mm}\textbf{Họ và tên:} \textbf{\large Đào Viết Tuấn Anh} \hfill \textbf{Mã HS:} \texttt{TUANANH2411} \quad \\
\hline
\end{tabularx}

\vspace{1.5mm}

\noindent\textbf{\large Câu 2 (2,5 điểm). Rút gọn và tìm $x$:}
\vspace{0.8mm}

\noindent
\begin{tabularx}{\linewidth}{|X|X|}
\hline
\oBaiTap[12]{a)}{Rút gọn và chứng minh giá trị biểu thức không phụ thuộc vào giá trị của biến $x$:\par
$M = (2x - 1)(x + 3) - 2x(x + 2) - x + 7$.} & 
\oBaiTap[12]{b)}{Tìm số thực $x$, biết:\par
$(3x - 1)(2x + 1) - 6x(x - 2) = 25$.} \\
\hline
\end{tabularx}

\vspace{1.8mm}

\noindent\textbf{\large Câu 3 (2,5 điểm). Bài toán thực tế và nâng cao:}
\vspace{0.8mm}

\noindent
\begin{tabularx}{\linewidth}{|X|}
\hline
\vspace{0.5mm}
\textbf{a) (1,5 điểm)} Một mảnh vườn hình chữ nhật có chiều dài là $(3x + 4)\text{ m}$ và chiều rộng là $(2x + 1)\text{ m}$. Người ta làm một lối đi xung quanh vườn rộng $1\text{ m}$, phần đất còn lại bên trong dùng để trồng hoa.\par
1) Viết đa thức theo biến $x$ biểu thị diện tích mảnh vườn ban đầu.\par
2) Viết đa thức theo biến $x$ biểu thị diện tích phần đất bên trong dùng để trồng hoa.\par
\vspace{0.4mm}
\noindent\textit{Bài giải:}\dongke{9}
\vspace{1mm}
\textbf{b) (1,0 điểm)} Cho hai số $x, y$ thỏa mãn $x + y = 4$ và $xy = 3$. Không tính cụ thể giá trị của $x$ và $y$, hãy tính giá trị của biểu thức $T = x^3 + y^3$.\par
\vspace{0.4mm}
\noindent\textit{Bài giải:}\dongke{7}
\\ \hline
\end{tabularx}


\newpage

% #############################################################################
% ========================= TOÁN 8 - TUẦN 03 - BUỔI 02: MÃ ĐỀ 102 ============
% ========================= MẶT TRƯỚC (TRANG 3) ===============================
% #############################################################################

\noindent
\begin{tabularx}{\linewidth}{|X|c|}
\hline
\textbf{PHIẾU KIỂM TRA ĐỊNH KỲ TOÁN 8 - TUẦN 03 - BUỔI 02} & \textbf{Mã đề: 102} \\
\textit{Chủ đề: Ôn tập Chương 1 - Đa thức} & \textit{Thời gian: 45 phút} \\
\hline
\end{tabularx}

\vspace{1mm}

\noindent
\begin{tabularx}{\linewidth}{|X|}
\hline
\rule[-2mm]{0pt}{7.5mm}\textbf{Họ và tên:} \textbf{\large Nguyễn Đình Gia Bảo} \hfill \textbf{Mã HS:} \texttt{GIABAO2612} \quad \\
\hline
\end{tabularx}

\vspace{1mm}

\noindent
\begin{tabularx}{\linewidth}{|p{3.2cm}|X|}
\hline
\centering\textbf{ĐIỂM SỐ} & \centering\textbf{LỜI NHẬN XÉT CỦA THẦY CÔ} \tabularnewline
\hline
\rule{0pt}{1.3cm} & \tabularnewline
\hline
\end{tabularx}

\vspace{1.2mm}

\noindent\textbf{BẢNG ĐIỀN ĐÁP ÁN TRẮC NGHIỆM (Mỗi câu đúng 0,5 điểm):}
\vspace{0.8mm}

\noindent
\renewcommand{\arraystretch}{1.35}
\begin{tabularx}{\linewidth}{|c|*{5}{>{\centering\arraybackslash}X|}}
\hline
\textbf{Câu} & \textbf{1} & \textbf{2} & \textbf{3} & \textbf{4} & \textbf{5} \\
\hline
\textbf{Đáp án} & & & & & \\
\hline
\end{tabularx}
\renewcommand{\arraystretch}{1.0}

\vspace{1.5mm}

\noindent\textbf{\large PHẦN I. TRẮC NGHIỆM KHÁCH QUAN (2,5 điểm)}
\vspace{1mm}

\cauTN{1}{Bậc của đa thức $N = 4x^3 y^3 - 3x^2 y^4 + 2x^5 y - 4x^3 y^3 + 9$ là:}{$5$}{$7$}{$6$}{$4$}

\cauTN{2}{Đơn thức nào sau đây đồng dạng với đơn thức $-5x^2 y^4$?}{$3x^2 y^4$}{$-5x^4 y^2$}{$2x^2 y^2$}{$4xy^4$}

\cauTN{3}{Kết quả của phép nhân đa thức $(2x + 3)(x - 2)$ là:}{$2x^2 + x - 6$}{$2x^2 - 7x - 6$}{$2x^2 - 6$}{$2x^2 - x - 6$}

\cauTN{4}{Kết quả của phép chia $(20x^5 y^3 - 15x^3 y^2) : 5x^3 y^2$ là:}{$4x^2 y - 3$}{$4x^2 y - 3xy$}{$4x^2 y^2 - 3$}{$4x^2 - 3y$}

\cauTN{5}{Giá trị của biểu thức $Q = x^2 + 6xy + 9y^2$ tại $x = 14$ và $y = 2$ là:}{$256$}{$400$}{$324$}{$100$}

\vspace{1mm}

\noindent\textbf{\large PHẦN II. TỰ LUẬN (7,5 điểm)}

\noindent\textbf{Câu 1 (2,5 điểm). Thực hiện phép tính với đa thức:}
\vspace{0.8mm}

\noindent
\begin{tabularx}{\linewidth}{|X|X|}
\hline
\multicolumn{2}{|p{0.97\linewidth}|}{Cho hai đa thức: $A = 4x^2 y - 3xy^2 + 4xy - 2$ \ và \ $B = 3x^2 y + 3xy^2 - 2xy + 5$.} \\
\hline
\oBaiTap[7]{a)}{Tính $C = A + B$:} & 
\oBaiTap[7]{b)}{Tính $D = A - B$:} \\
\hline
\oBaiTap[5]{c)}{Tìm bậc của đa thức $C$:} & 
\oBaiTap[5]{d)}{Tính giá trị của đa thức $C$ tại $x = 1; y = -1$:} \\
\hline
\end{tabularx}


\newpage

% #############################################################################
% ========================= TOÁN 8 - TUẦN 03 - BUỔI 02: MÃ ĐỀ 102 ============
% ========================= MẶT SAU (TRANG 4) =================================
% #############################################################################

\noindent
\begin{tabularx}{\linewidth}{|X|c|}
\hline
\textbf{PHẦN II. TỰ LUẬN (Tiếp theo) - TOÁN 8 TUẦN 03 BUỔI 02} & \textbf{Mã đề: 102} \\
\hline
\end{tabularx}

\vspace{1mm}

\noindent
\begin{tabularx}{\linewidth}{|X|}
\hline
\rule[-2mm]{0pt}{7mm}\textbf{Họ và tên:} \textbf{\large Nguyễn Đình Gia Bảo} \hfill \textbf{Mã HS:} \texttt{GIABAO2612} \quad \\
\hline
\end{tabularx}

\vspace{1.5mm}

\noindent\textbf{\large Câu 2 (2,5 điểm). Rút gọn và tìm $x$:}
\vspace{0.8mm}

\noindent
\begin{tabularx}{\linewidth}{|X|X|}
\hline
\oBaiTap[12]{a)}{Rút gọn và chứng minh giá trị biểu thức không phụ thuộc vào giá trị của biến $x$:\par
$N = (3x - 1)(x + 2) - 3x(x + 1) - 2x + 5$.} & 
\oBaiTap[12]{b)}{Tìm số thực $x$, biết:\par
$(2x - 1)(3x + 2) - 6x(x - 1) = 19$.} \\
\hline
\end{tabularx}

\vspace{1.8mm}

\noindent\textbf{\large Câu 3 (2,5 điểm). Bài toán thực tế và nâng cao:}
\vspace{0.8mm}

\noindent
\begin{tabularx}{\linewidth}{|X|}
\hline
\vspace{0.5mm}
\textbf{a) (1,5 điểm)} Một khu đất hình chữ nhật có chiều dài là $(4x + 3)\text{ m}$ và chiều rộng là $(2x + 3)\text{ m}$. Người ta làm một lối đi xung quanh khu đất rộng $1\text{ m}$, phần đất còn lại bên trong dùng để trồng rau sạch.\par
1) Viết đa thức theo biến $x$ biểu thị diện tích khu đất ban đầu.\par
2) Viết đa thức theo biến $x$ biểu thị diện tích phần đất bên trong dùng để trồng rau.\par
\vspace{0.4mm}
\noindent\textit{Bài giải:}\dongke{9}
\vspace{1mm}
\textbf{b) (1,0 điểm)} Cho hai số $x, y$ thỏa mãn $x + y = 5$ và $xy = 4$. Không tính cụ thể giá trị của $x$ và $y$, hãy tính giá trị của biểu thức $M = x^3 + y^3$.\par
\vspace{0.4mm}
\noindent\textit{Bài giải:}\dongke{7}
\\ \hline
\end{tabularx}


\newpage

% #############################################################################
% ========================= TOÁN 8 - TUẦN 03 - BUỔI 02: MÃ ĐỀ 103 ============
% ========================= MẶT TRƯỚC (TRANG 5) ===============================
% #############################################################################

\noindent
\begin{tabularx}{\linewidth}{|X|c|}
\hline
\textbf{PHIẾU KIỂM TRA ĐỊNH KỲ TOÁN 8 - TUẦN 03 - BUỔI 02} & \textbf{Mã đề: 103} \\
\textit{Chủ đề: Ôn tập Chương 1 - Đa thức} & \textit{Thời gian: 45 phút} \\
\hline
\end{tabularx}

\vspace{1mm}

\noindent
\begin{tabularx}{\linewidth}{|X|}
\hline
\rule[-2mm]{0pt}{7.5mm}\textbf{Họ và tên:} \textbf{\large Phan Gia Bảo} \hfill \textbf{Mã HS:} \texttt{GIABAO1804} \quad \\
\hline
\end{tabularx}

\vspace{1mm}

\noindent
\begin{tabularx}{\linewidth}{|p{3.2cm}|X|}
\hline
\centering\textbf{ĐIỂM SỐ} & \centering\textbf{LỜI NHẬN XÉT CỦA THẦY CÔ} \tabularnewline
\hline
\rule{0pt}{1.3cm} & \tabularnewline
\hline
\end{tabularx}

\vspace{1.2mm}

\noindent\textbf{BẢNG ĐIỀN ĐÁP ÁN TRẮC NGHIỆM (Mỗi câu đúng 0,5 điểm):}
\vspace{0.8mm}

\noindent
\renewcommand{\arraystretch}{1.35}
\begin{tabularx}{\linewidth}{|c|*{5}{>{\centering\arraybackslash}X|}}
\hline
\textbf{Câu} & \textbf{1} & \textbf{2} & \textbf{3} & \textbf{4} & \textbf{5} \\
\hline
\textbf{Đáp án} & & & & & \\
\hline
\end{tabularx}
\renewcommand{\arraystretch}{1.0}

\vspace{1.5mm}

\noindent\textbf{\large PHẦN I. TRẮC NGHIỆM KHÁCH QUAN (2,5 điểm)}
\vspace{1mm}

\cauTN{1}{Bậc của đa thức $P = 6x^4 y^2 - 3x^3 y^3 + 2x^5 y - 6x^4 y^2 + 8$ là:}{$5$}{$6$}{$7$}{$4$}

\cauTN{2}{Đơn thức nào sau đây đồng dạng với đơn thức $-6x^3 y^3$?}{$5x^3 y^3$}{$-6x^2 y^3$}{$3x^3 y^2$}{$4x^4 y^2$}

\cauTN{3}{Kết quả của phép nhân đa thức $(3x - 2)(x + 2)$ là:}{$3x^2 + 4x - 4$}{$3x^2 - 4x - 4$}{$3x^2 + 6x - 4$}{$3x^2 - 4$}

\cauTN{4}{Kết quả của phép chia $(24x^4 y^3 - 16x^3 y^2) : 8x^2 y^2$ là:}{$3x^2 y - 2x$}{$3x^2 y - 2$}{$3x^2 y^2 - 2x$}{$3x^2 - 2y$}

\cauTN{5}{Giá trị của biểu thức $M = x^2 - 8xy + 16y^2$ tại $x = 18$ và $y = 2$ là:}{$64$}{$144$}{$100$}{$400$}

\vspace{1mm}

\noindent\textbf{\large PHẦN II. TỰ LUẬN (7,5 điểm)}

\noindent\textbf{Câu 1 (2,5 điểm). Thực hiện phép tính với đa thức:}
\vspace{0.8mm}

\noindent
\begin{tabularx}{\linewidth}{|X|X|}
\hline
\multicolumn{2}{|p{0.97\linewidth}|}{Cho hai đa thức: $A = 5x^2 y - 3xy^2 + 6xy - 3$ \ và \ $B = 2x^2 y + 3xy^2 - 4xy + 5$.} \\
\hline
\oBaiTap[7]{a)}{Tính $C = A + B$:} & 
\oBaiTap[7]{b)}{Tính $D = A - B$:} \\
\hline
\oBaiTap[5]{c)}{Tìm bậc của đa thức $C$:} & 
\oBaiTap[5]{d)}{Tính giá trị của đa thức $C$ tại $x = 1; y = -1$:} \\
\hline
\end{tabularx}


\newpage

% #############################################################################
% ========================= TOÁN 8 - TUẦN 03 - BUỔI 02: MÃ ĐỀ 103 ============
% ========================= MẶT SAU (TRANG 6) =================================
% #############################################################################

\noindent
\begin{tabularx}{\linewidth}{|X|c|}
\hline
\textbf{PHẦN II. TỰ LUẬN (Tiếp theo) - TOÁN 8 TUẦN 03 BUỔI 02} & \textbf{Mã đề: 103} \\
\hline
\end{tabularx}

\vspace{1mm}

\noindent
\begin{tabularx}{\linewidth}{|X|}
\hline
\rule[-2mm]{0pt}{7mm}\textbf{Họ và tên:} \textbf{\large Phan Gia Bảo} \hfill \textbf{Mã HS:} \texttt{GIABAO1804} \quad \\
\hline
\end{tabularx}

\vspace{1.5mm}

\noindent\textbf{\large Câu 2 (2,5 điểm). Rút gọn và tìm $x$:}
\vspace{0.8mm}

\noindent
\begin{tabularx}{\linewidth}{|X|X|}
\hline
\oBaiTap[12]{a)}{Rút gọn và chứng minh giá trị biểu thức không phụ thuộc vào giá trị của biến $x$:\par
$P = (2x - 3)(x + 2) - 2x(x + 1) + x + 8$.} & 
\oBaiTap[12]{b)}{Tìm số thực $x$, biết:\par
$(3x - 2)(2x + 1) - 6x(x - 1) = 23$.} \\
\hline
\end{tabularx}

\vspace{1.8mm}

\noindent\textbf{\large Câu 3 (2,5 điểm). Bài toán thực tế và nâng cao:}
\vspace{0.8mm}

\noindent
\begin{tabularx}{\linewidth}{|X|}
\hline
\vspace{0.5mm}
\textbf{a) (1,5 điểm)} Một mảnh vườn hình chữ nhật có chiều dài là $(3x + 5)\text{ m}$ và chiều rộng là $(2x + 2)\text{ m}$. Người ta làm một lối đi xung quanh vườn rộng $1\text{ m}$, phần đất còn lại bên trong dùng để trồng hoa.\par
1) Viết đa thức theo biến $x$ biểu thị diện tích mảnh vườn ban đầu.\par
2) Viết đa thức theo biến $x$ biểu thị diện tích phần đất bên trong dùng để trồng hoa.\par
\vspace{0.4mm}
\noindent\textit{Bài giải:}\dongke{9}
\vspace{1mm}
\textbf{b) (1,0 điểm)} Cho hai số $x, y$ thỏa mãn $x + y = 6$ và $xy = 5$. Không tính cụ thể giá trị của $x$ và $y$, hãy tính giá trị của biểu thức $T = x^3 + y^3$.\par
\vspace{0.4mm}
\noindent\textit{Bài giải:}\dongke{7}
\\ \hline
\end{tabularx}


\newpage

% #############################################################################
% ========================= TOÁN 8 - TUẦN 03 - BUỔI 02: MÃ ĐỀ 104 ============
% ========================= MẶT TRƯỚC (TRANG 7) ===============================
% #############################################################################

\noindent
\begin{tabularx}{\linewidth}{|X|c|}
\hline
\textbf{PHIẾU KIỂM TRA ĐỊNH KỲ TOÁN 8 - TUẦN 03 - BUỔI 02} & \textbf{Mã đề: 104} \\
\textit{Chủ đề: Ôn tập Chương 1 - Đa thức} & \textit{Thời gian: 45 phút} \\
\hline
\end{tabularx}

\vspace{1mm}

\noindent
\begin{tabularx}{\linewidth}{|X|}
\hline
\rule[-2mm]{0pt}{7.5mm}\textbf{Họ và tên:} \textbf{\large Nguyễn Thùy Dung} \hfill \textbf{Mã HS:} \texttt{THUYDUNG1208} \quad \\
\hline
\end{tabularx}

\vspace{1mm}

\noindent
\begin{tabularx}{\linewidth}{|p{3.2cm}|X|}
\hline
\centering\textbf{ĐIỂM SỐ} & \centering\textbf{LỜI NHẬN XÉT CỦA THẦY CÔ} \tabularnewline
\hline
\rule{0pt}{1.3cm} & \tabularnewline
\hline
\end{tabularx}

\vspace{1.2mm}

\noindent\textbf{BẢNG ĐIỀN ĐÁP ÁN TRẮC NGHIỆM (Mỗi câu đúng 0,5 điểm):}
\vspace{0.8mm}

\noindent
\renewcommand{\arraystretch}{1.35}
\begin{tabularx}{\linewidth}{|c|*{5}{>{\centering\arraybackslash}X|}}
\hline
\textbf{Câu} & \textbf{1} & \textbf{2} & \textbf{3} & \textbf{4} & \textbf{5} \\
\hline
\textbf{Đáp án} & & & & & \\
\hline
\end{tabularx}
\renewcommand{\arraystretch}{1.0}

\vspace{1.5mm}

\noindent\textbf{\large PHẦN I. TRẮC NGHIỆM KHÁCH QUAN (2,5 điểm)}
\vspace{1mm}

\cauTN{1}{Bậc của đa thức $Q = 5x^4 y^3 - 2x^3 y^4 + 4x^6 y - 5x^4 y^3 + 7$ là:}{$6$}{$5$}{$7$}{$4$}

\cauTN{2}{Đơn thức nào sau đây đồng dạng với đơn thức $-7x^4 y^2$?}{$3x^4 y^2$}{$-7x^2 y^4$}{$4x^4 y^3$}{$2x^3 y^2$}

\cauTN{3}{Kết quả của phép nhân đa thức $(3x + 2)(x - 2)$ là:}{$3x^2 + 4x - 4$}{$3x^2 - 4x - 4$}{$3x^2 - 6$}{$3x^2 - 8x - 4$}

\cauTN{4}{Kết quả của phép chia $(24x^5 y^3 - 18x^3 y^2) : 6x^3 y^2$ là:}{$4x^2 y - 3$}{$4x^2 y - 3xy$}{$4x^2 y^2 - 3$}{$4x^2 - 3y$}

\cauTN{5}{Giá trị của biểu thức $N = x^2 + 8xy + 16y^2$ tại $x = 12$ và $y = 2$ là:}{$256$}{$400$}{$324$}{$100$}

\vspace{1mm}

\noindent\textbf{\large PHẦN II. TỰ LUẬN (7,5 điểm)}

\noindent\textbf{Câu 1 (2,5 điểm). Thực hiện phép tính với đa thức:}
\vspace{0.8mm}

\noindent
\begin{tabularx}{\linewidth}{|X|X|}
\hline
\multicolumn{2}{|p{0.97\linewidth}|}{Cho hai đa thức: $A = 5x^2 y - 4xy^2 + 3xy - 1$ \ và \ $B = 3x^2 y + 4xy^2 - 2xy + 6$.} \\
\hline
\oBaiTap[7]{a)}{Tính $C = A + B$:} & 
\oBaiTap[7]{b)}{Tính $D = A - B$:} \\
\hline
\oBaiTap[5]{c)}{Tìm bậc của đa thức $C$:} & 
\oBaiTap[5]{d)}{Tính giá trị của đa thức $C$ tại $x = 1; y = -1$:} \\
\hline
\end{tabularx}


\newpage

% #############################################################################
% ========================= TOÁN 8 - TUẦN 03 - BUỔI 02: MÃ ĐỀ 104 ============
% ========================= MẶT SAU (TRANG 8) =================================
% #############################################################################

\noindent
\begin{tabularx}{\linewidth}{|X|c|}
\hline
\textbf{PHẦN II. TỰ LUẬN (Tiếp theo) - TOÁN 8 TUẦN 03 BUỔI 02} & \textbf{Mã đề: 104} \\
\hline
\end{tabularx}

\vspace{1mm}

\noindent
\begin{tabularx}{\linewidth}{|X|}
\hline
\rule[-2mm]{0pt}{7mm}\textbf{Họ và tên:} \textbf{\large Nguyễn Thùy Dung} \hfill \textbf{Mã HS:} \texttt{THUYDUNG1208} \quad \\
\hline
\end{tabularx}

\vspace{1.5mm}

\noindent\textbf{\large Câu 2 (2,5 điểm). Rút gọn và tìm $x$:}
\vspace{0.8mm}

\noindent
\begin{tabularx}{\linewidth}{|X|X|}
\hline
\oBaiTap[12]{a)}{Rút gọn và chứng minh giá trị biểu thức không phụ thuộc vào giá trị của biến $x$:\par
$N = (4x - 1)(x + 2) - 4x(x + 1) - 3x + 9$.} & 
\oBaiTap[12]{b)}{Tìm số thực $x$, biết:\par
$(2x - 1)(4x + 3) - 8x(x - 1) = 29$.} \\
\hline
\end{tabularx}

\vspace{1.8mm}

\noindent\textbf{\large Câu 3 (2,5 điểm). Bài toán thực tế và nâng cao:}
\vspace{0.8mm}

\noindent
\begin{tabularx}{\linewidth}{|X|}
\hline
\vspace{0.5mm}
\textbf{a) (1,5 điểm)} Một khu đất hình chữ nhật có chiều dài là $(4x + 5)\text{ m}$ và chiều rộng là $(2x + 4)\text{ m}$. Người ta làm một lối đi xung quanh khu đất rộng $1\text{ m}$, phần đất còn lại bên trong dùng để trồng rau sạch.\par
1) Viết đa thức theo biến $x$ biểu thị diện tích khu đất ban đầu.\par
2) Viết đa thức theo biến $x$ biểu thị diện tích phần đất bên trong dùng để trồng rau.\par
\vspace{0.4mm}
\noindent\textit{Bài giải:}\dongke{9}
\vspace{1mm}
\textbf{b) (1,0 điểm)} Cho hai số $x, y$ thỏa mãn $x + y = 6$ và $xy = 8$. Không tính cụ thể giá trị của $x$ và $y$, hãy tính giá trị của biểu thức $M = x^3 + y^3$.\par
\vspace{0.4mm}
\noindent\textit{Bài giải:}\dongke{7}
\\ \hline
\end{tabularx}



\newpage


% #############################################################################
% ========================= TOÁN 8 - TUẦN 03 - BUỔI 02: MÃ ĐỀ 101 ============
% ========================= MẶT TRƯỚC (TRANG 1) ===============================
% #############################################################################

\noindent
\begin{tabularx}{\linewidth}{|X|c|}
\hline
\textbf{PHIẾU KIỂM TRA ĐỊNH KỲ TOÁN 8 - TUẦN 03 - BUỔI 02} & \textbf{Mã đề: 101} \\
\textit{Chủ đề: Ôn tập Chương 1 - Đa thức} & \textit{Thời gian: 45 phút} \\
\hline
\end{tabularx}

\vspace{1mm}

\noindent
\begin{tabularx}{\linewidth}{|X|}
\hline
\rule[-2mm]{0pt}{7.5mm}\textbf{Họ và tên:} \textbf{\large Lê Ngọc Gia Hân} \hfill \textbf{Mã HS:} \texttt{GIAHAN2211} \quad \\
\hline
\end{tabularx}

\vspace{1mm}

\noindent
\begin{tabularx}{\linewidth}{|p{3.2cm}|X|}
\hline
\centering\textbf{ĐIỂM SỐ} & \centering\textbf{LỜI NHẬN XÉT CỦA THẦY CÔ} \tabularnewline
\hline
\rule{0pt}{1.3cm} & \tabularnewline
\hline
\end{tabularx}

\vspace{1.2mm}

\noindent\textbf{BẢNG ĐIỀN ĐÁP ÁN TRẮC NGHIỆM (Mỗi câu đúng 0,5 điểm):}
\vspace{0.8mm}

\noindent
\renewcommand{\arraystretch}{1.35}
\begin{tabularx}{\linewidth}{|c|*{5}{>{\centering\arraybackslash}X|}}
\hline
\textbf{Câu} & \textbf{1} & \textbf{2} & \textbf{3} & \textbf{4} & \textbf{5} \\
\hline
\textbf{Đáp án} & & & & & \\
\hline
\end{tabularx}
\renewcommand{\arraystretch}{1.0}

\vspace{1.5mm}

\noindent\textbf{\large PHẦN I. TRẮC NGHIỆM KHÁCH QUAN (2,5 điểm)}
\vspace{1mm}

\cauTN{1}{Bậc của đa thức $M = 5x^3 y^2 - 2x^2 y^3 + 3x^4 y - 5x^3 y^2 + 1$ là:}{$4$}{$5$}{$6$}{$3$}

\cauTN{2}{Đơn thức nào sau đây đồng dạng với đơn thức $-4x^3 y^2$?}{$2x^2 y^3$}{$-4x^3 y$}{$7x^3 y^2$}{$5x^2 y^2$}

\cauTN{3}{Kết quả của phép nhân đa thức $(2x - 3)(x + 2)$ là:}{$2x^2 + x - 6$}{$2x^2 - x - 6$}{$2x^2 + 7x - 6$}{$2x^2 - 6$}

\cauTN{4}{Kết quả của phép chia $(18x^4 y^3 - 12x^3 y^2) : 6x^2 y^2$ là:}{$3x^2 y - 2xy$}{$3x^2 y - 2$}{$3x^2 y^2 - 2x$}{$3x^2 y - 2x$}

\cauTN{5}{Giá trị của biểu thức $P = x^2 - 6xy + 9y^2$ tại $x = 16$ và $y = 2$ là:}{$100$}{$144$}{$64$}{$400$}

\vspace{1mm}

\noindent\textbf{\large PHẦN II. TỰ LUẬN (7,5 điểm)}

\noindent\textbf{Câu 1 (2,5 điểm). Thực hiện phép tính với đa thức:}
\vspace{0.8mm}

\noindent
\begin{tabularx}{\linewidth}{|X|X|}
\hline
\multicolumn{2}{|p{0.97\linewidth}|}{Cho hai đa thức: $A = 3x^2 y - 2xy^2 + 5xy - 1$ \ và \ $B = 2x^2 y + 2xy^2 - 3xy + 4$.} \\
\hline
\oBaiTap[7]{a)}{Tính $C = A + B$:} & 
\oBaiTap[7]{b)}{Tính $D = A - B$:} \\
\hline
\oBaiTap[5]{c)}{Tìm bậc của đa thức $C$:} & 
\oBaiTap[5]{d)}{Tính giá trị của đa thức $C$ tại $x = 1; y = -1$:} \\
\hline
\end{tabularx}


\newpage

% #############################################################################
% ========================= TOÁN 8 - TUẦN 03 - BUỔI 02: MÃ ĐỀ 101 ============
% ========================= MẶT SAU (TRANG 2) =================================
% #############################################################################

\noindent
\begin{tabularx}{\linewidth}{|X|c|}
\hline
\textbf{PHẦN II. TỰ LUẬN (Tiếp theo) - TOÁN 8 TUẦN 03 BUỔI 02} & \textbf{Mã đề: 101} \\
\hline
\end{tabularx}

\vspace{1mm}

\noindent
\begin{tabularx}{\linewidth}{|X|}
\hline
\rule[-2mm]{0pt}{7mm}\textbf{Họ và tên:} \textbf{\large Lê Ngọc Gia Hân} \hfill \textbf{Mã HS:} \texttt{GIAHAN2211} \quad \\
\hline
\end{tabularx}

\vspace{1.5mm}

\noindent\textbf{\large Câu 2 (2,5 điểm). Rút gọn và tìm $x$:}
\vspace{0.8mm}

\noindent
\begin{tabularx}{\linewidth}{|X|X|}
\hline
\oBaiTap[12]{a)}{Rút gọn và chứng minh giá trị biểu thức không phụ thuộc vào giá trị của biến $x$:\par
$M = (2x - 1)(x + 3) - 2x(x + 2) - x + 7$.} & 
\oBaiTap[12]{b)}{Tìm số thực $x$, biết:\par
$(3x - 1)(2x + 1) - 6x(x - 2) = 25$.} \\
\hline
\end{tabularx}

\vspace{1.8mm}

\noindent\textbf{\large Câu 3 (2,5 điểm). Bài toán thực tế và nâng cao:}
\vspace{0.8mm}

\noindent
\begin{tabularx}{\linewidth}{|X|}
\hline
\vspace{0.5mm}
\textbf{a) (1,5 điểm)} Một mảnh vườn hình chữ nhật có chiều dài là $(3x + 4)\text{ m}$ và chiều rộng là $(2x + 1)\text{ m}$. Người ta làm một lối đi xung quanh vườn rộng $1\text{ m}$, phần đất còn lại bên trong dùng để trồng hoa.\par
1) Viết đa thức theo biến $x$ biểu thị diện tích mảnh vườn ban đầu.\par
2) Viết đa thức theo biến $x$ biểu thị diện tích phần đất bên trong dùng để trồng hoa.\par
\vspace{0.4mm}
\noindent\textit{Bài giải:}\dongke{9}
\vspace{1mm}
\textbf{b) (1,0 điểm)} Cho hai số $x, y$ thỏa mãn $x + y = 4$ và $xy = 3$. Không tính cụ thể giá trị của $x$ và $y$, hãy tính giá trị của biểu thức $T = x^3 + y^3$.\par
\vspace{0.4mm}
\noindent\textit{Bài giải:}\dongke{7}
\\ \hline
\end{tabularx}


\newpage

% #############################################################################
% ========================= TOÁN 8 - TUẦN 03 - BUỔI 02: MÃ ĐỀ 102 ============
% ========================= MẶT TRƯỚC (TRANG 3) ===============================
% #############################################################################

\noindent
\begin{tabularx}{\linewidth}{|X|c|}
\hline
\textbf{PHIẾU KIỂM TRA ĐỊNH KỲ TOÁN 8 - TUẦN 03 - BUỔI 02} & \textbf{Mã đề: 102} \\
\textit{Chủ đề: Ôn tập Chương 1 - Đa thức} & \textit{Thời gian: 45 phút} \\
\hline
\end{tabularx}

\vspace{1mm}

\noindent
\begin{tabularx}{\linewidth}{|X|}
\hline
\rule[-2mm]{0pt}{7.5mm}\textbf{Họ và tên:} \textbf{\large Trần Đăng Khôi} \hfill \textbf{Mã HS:} \texttt{DANGKHOI2510} \quad \\
\hline
\end{tabularx}

\vspace{1mm}

\noindent
\begin{tabularx}{\linewidth}{|p{3.2cm}|X|}
\hline
\centering\textbf{ĐIỂM SỐ} & \centering\textbf{LỜI NHẬN XÉT CỦA THẦY CÔ} \tabularnewline
\hline
\rule{0pt}{1.3cm} & \tabularnewline
\hline
\end{tabularx}

\vspace{1.2mm}

\noindent\textbf{BẢNG ĐIỀN ĐÁP ÁN TRẮC NGHIỆM (Mỗi câu đúng 0,5 điểm):}
\vspace{0.8mm}

\noindent
\renewcommand{\arraystretch}{1.35}
\begin{tabularx}{\linewidth}{|c|*{5}{>{\centering\arraybackslash}X|}}
\hline
\textbf{Câu} & \textbf{1} & \textbf{2} & \textbf{3} & \textbf{4} & \textbf{5} \\
\hline
\textbf{Đáp án} & & & & & \\
\hline
\end{tabularx}
\renewcommand{\arraystretch}{1.0}

\vspace{1.5mm}

\noindent\textbf{\large PHẦN I. TRẮC NGHIỆM KHÁCH QUAN (2,5 điểm)}
\vspace{1mm}

\cauTN{1}{Bậc của đa thức $N = 4x^3 y^3 - 3x^2 y^4 + 2x^5 y - 4x^3 y^3 + 9$ là:}{$5$}{$7$}{$6$}{$4$}

\cauTN{2}{Đơn thức nào sau đây đồng dạng với đơn thức $-5x^2 y^4$?}{$3x^2 y^4$}{$-5x^4 y^2$}{$2x^2 y^2$}{$4xy^4$}

\cauTN{3}{Kết quả của phép nhân đa thức $(2x + 3)(x - 2)$ là:}{$2x^2 + x - 6$}{$2x^2 - 7x - 6$}{$2x^2 - 6$}{$2x^2 - x - 6$}

\cauTN{4}{Kết quả của phép chia $(20x^5 y^3 - 15x^3 y^2) : 5x^3 y^2$ là:}{$4x^2 y - 3$}{$4x^2 y - 3xy$}{$4x^2 y^2 - 3$}{$4x^2 - 3y$}

\cauTN{5}{Giá trị của biểu thức $Q = x^2 + 6xy + 9y^2$ tại $x = 14$ và $y = 2$ là:}{$256$}{$400$}{$324$}{$100$}

\vspace{1mm}

\noindent\textbf{\large PHẦN II. TỰ LUẬN (7,5 điểm)}

\noindent\textbf{Câu 1 (2,5 điểm). Thực hiện phép tính với đa thức:}
\vspace{0.8mm}

\noindent
\begin{tabularx}{\linewidth}{|X|X|}
\hline
\multicolumn{2}{|p{0.97\linewidth}|}{Cho hai đa thức: $A = 4x^2 y - 3xy^2 + 4xy - 2$ \ và \ $B = 3x^2 y + 3xy^2 - 2xy + 5$.} \\
\hline
\oBaiTap[7]{a)}{Tính $C = A + B$:} & 
\oBaiTap[7]{b)}{Tính $D = A - B$:} \\
\hline
\oBaiTap[5]{c)}{Tìm bậc của đa thức $C$:} & 
\oBaiTap[5]{d)}{Tính giá trị của đa thức $C$ tại $x = 1; y = -1$:} \\
\hline
\end{tabularx}


\newpage

% #############################################################################
% ========================= TOÁN 8 - TUẦN 03 - BUỔI 02: MÃ ĐỀ 102 ============
% ========================= MẶT SAU (TRANG 4) =================================
% #############################################################################

\noindent
\begin{tabularx}{\linewidth}{|X|c|}
\hline
\textbf{PHẦN II. TỰ LUẬN (Tiếp theo) - TOÁN 8 TUẦN 03 BUỔI 02} & \textbf{Mã đề: 102} \\
\hline
\end{tabularx}

\vspace{1mm}

\noindent
\begin{tabularx}{\linewidth}{|X|}
\hline
\rule[-2mm]{0pt}{7mm}\textbf{Họ và tên:} \textbf{\large Trần Đăng Khôi} \hfill \textbf{Mã HS:} \texttt{DANGKHOI2510} \quad \\
\hline
\end{tabularx}

\vspace{1.5mm}

\noindent\textbf{\large Câu 2 (2,5 điểm). Rút gọn và tìm $x$:}
\vspace{0.8mm}

\noindent
\begin{tabularx}{\linewidth}{|X|X|}
\hline
\oBaiTap[12]{a)}{Rút gọn và chứng minh giá trị biểu thức không phụ thuộc vào giá trị của biến $x$:\par
$N = (3x - 1)(x + 2) - 3x(x + 1) - 2x + 5$.} & 
\oBaiTap[12]{b)}{Tìm số thực $x$, biết:\par
$(2x - 1)(3x + 2) - 6x(x - 1) = 19$.} \\
\hline
\end{tabularx}

\vspace{1.8mm}

\noindent\textbf{\large Câu 3 (2,5 điểm). Bài toán thực tế và nâng cao:}
\vspace{0.8mm}

\noindent
\begin{tabularx}{\linewidth}{|X|}
\hline
\vspace{0.5mm}
\textbf{a) (1,5 điểm)} Một khu đất hình chữ nhật có chiều dài là $(4x + 3)\text{ m}$ và chiều rộng là $(2x + 3)\text{ m}$. Người ta làm một lối đi xung quanh khu đất rộng $1\text{ m}$, phần đất còn lại bên trong dùng để trồng rau sạch.\par
1) Viết đa thức theo biến $x$ biểu thị diện tích khu đất ban đầu.\par
2) Viết đa thức theo biến $x$ biểu thị diện tích phần đất bên trong dùng để trồng rau.\par
\vspace{0.4mm}
\noindent\textit{Bài giải:}\dongke{9}
\vspace{1mm}
\textbf{b) (1,0 điểm)} Cho hai số $x, y$ thỏa mãn $x + y = 5$ và $xy = 4$. Không tính cụ thể giá trị của $x$ và $y$, hãy tính giá trị của biểu thức $M = x^3 + y^3$.\par
\vspace{0.4mm}
\noindent\textit{Bài giải:}\dongke{7}
\\ \hline
\end{tabularx}


\newpage

% #############################################################################
% ========================= TOÁN 8 - TUẦN 03 - BUỔI 02: MÃ ĐỀ 103 ============
% ========================= MẶT TRƯỚC (TRANG 5) ===============================
% #############################################################################

\noindent
\begin{tabularx}{\linewidth}{|X|c|}
\hline
\textbf{PHIẾU KIỂM TRA ĐỊNH KỲ TOÁN 8 - TUẦN 03 - BUỔI 02} & \textbf{Mã đề: 103} \\
\textit{Chủ đề: Ôn tập Chương 1 - Đa thức} & \textit{Thời gian: 45 phút} \\
\hline
\end{tabularx}

\vspace{1mm}

\noindent
\begin{tabularx}{\linewidth}{|X|}
\hline
\rule[-2mm]{0pt}{7.5mm}\textbf{Họ và tên:} \textbf{\large Nguyễn Khánh Linh} \hfill \textbf{Mã HS:} \texttt{KHANHLINH2211} \quad \\
\hline
\end{tabularx}

\vspace{1mm}

\noindent
\begin{tabularx}{\linewidth}{|p{3.2cm}|X|}
\hline
\centering\textbf{ĐIỂM SỐ} & \centering\textbf{LỜI NHẬN XÉT CỦA THẦY CÔ} \tabularnewline
\hline
\rule{0pt}{1.3cm} & \tabularnewline
\hline
\end{tabularx}

\vspace{1.2mm}

\noindent\textbf{BẢNG ĐIỀN ĐÁP ÁN TRẮC NGHIỆM (Mỗi câu đúng 0,5 điểm):}
\vspace{0.8mm}

\noindent
\renewcommand{\arraystretch}{1.35}
\begin{tabularx}{\linewidth}{|c|*{5}{>{\centering\arraybackslash}X|}}
\hline
\textbf{Câu} & \textbf{1} & \textbf{2} & \textbf{3} & \textbf{4} & \textbf{5} \\
\hline
\textbf{Đáp án} & & & & & \\
\hline
\end{tabularx}
\renewcommand{\arraystretch}{1.0}

\vspace{1.5mm}

\noindent\textbf{\large PHẦN I. TRẮC NGHIỆM KHÁCH QUAN (2,5 điểm)}
\vspace{1mm}

\cauTN{1}{Bậc của đa thức $P = 6x^4 y^2 - 3x^3 y^3 + 2x^5 y - 6x^4 y^2 + 8$ là:}{$5$}{$6$}{$7$}{$4$}

\cauTN{2}{Đơn thức nào sau đây đồng dạng với đơn thức $-6x^3 y^3$?}{$5x^3 y^3$}{$-6x^2 y^3$}{$3x^3 y^2$}{$4x^4 y^2$}

\cauTN{3}{Kết quả của phép nhân đa thức $(3x - 2)(x + 2)$ là:}{$3x^2 + 4x - 4$}{$3x^2 - 4x - 4$}{$3x^2 + 6x - 4$}{$3x^2 - 4$}

\cauTN{4}{Kết quả của phép chia $(24x^4 y^3 - 16x^3 y^2) : 8x^2 y^2$ là:}{$3x^2 y - 2x$}{$3x^2 y - 2$}{$3x^2 y^2 - 2x$}{$3x^2 - 2y$}

\cauTN{5}{Giá trị của biểu thức $M = x^2 - 8xy + 16y^2$ tại $x = 18$ và $y = 2$ là:}{$64$}{$144$}{$100$}{$400$}

\vspace{1mm}

\noindent\textbf{\large PHẦN II. TỰ LUẬN (7,5 điểm)}

\noindent\textbf{Câu 1 (2,5 điểm). Thực hiện phép tính với đa thức:}
\vspace{0.8mm}

\noindent
\begin{tabularx}{\linewidth}{|X|X|}
\hline
\multicolumn{2}{|p{0.97\linewidth}|}{Cho hai đa thức: $A = 5x^2 y - 3xy^2 + 6xy - 3$ \ và \ $B = 2x^2 y + 3xy^2 - 4xy + 5$.} \\
\hline
\oBaiTap[7]{a)}{Tính $C = A + B$:} & 
\oBaiTap[7]{b)}{Tính $D = A - B$:} \\
\hline
\oBaiTap[5]{c)}{Tìm bậc của đa thức $C$:} & 
\oBaiTap[5]{d)}{Tính giá trị của đa thức $C$ tại $x = 1; y = -1$:} \\
\hline
\end{tabularx}


\newpage

% #############################################################################
% ========================= TOÁN 8 - TUẦN 03 - BUỔI 02: MÃ ĐỀ 103 ============
% ========================= MẶT SAU (TRANG 6) =================================
% #############################################################################

\noindent
\begin{tabularx}{\linewidth}{|X|c|}
\hline
\textbf{PHẦN II. TỰ LUẬN (Tiếp theo) - TOÁN 8 TUẦN 03 BUỔI 02} & \textbf{Mã đề: 103} \\
\hline
\end{tabularx}

\vspace{1mm}

\noindent
\begin{tabularx}{\linewidth}{|X|}
\hline
\rule[-2mm]{0pt}{7mm}\textbf{Họ và tên:} \textbf{\large Nguyễn Khánh Linh} \hfill \textbf{Mã HS:} \texttt{KHANHLINH2211} \quad \\
\hline
\end{tabularx}

\vspace{1.5mm}

\noindent\textbf{\large Câu 2 (2,5 điểm). Rút gọn và tìm $x$:}
\vspace{0.8mm}

\noindent
\begin{tabularx}{\linewidth}{|X|X|}
\hline
\oBaiTap[12]{a)}{Rút gọn và chứng minh giá trị biểu thức không phụ thuộc vào giá trị của biến $x$:\par
$P = (2x - 3)(x + 2) - 2x(x + 1) + x + 8$.} & 
\oBaiTap[12]{b)}{Tìm số thực $x$, biết:\par
$(3x - 2)(2x + 1) - 6x(x - 1) = 23$.} \\
\hline
\end{tabularx}

\vspace{1.8mm}

\noindent\textbf{\large Câu 3 (2,5 điểm). Bài toán thực tế và nâng cao:}
\vspace{0.8mm}

\noindent
\begin{tabularx}{\linewidth}{|X|}
\hline
\vspace{0.5mm}
\textbf{a) (1,5 điểm)} Một mảnh vườn hình chữ nhật có chiều dài là $(3x + 5)\text{ m}$ và chiều rộng là $(2x + 2)\text{ m}$. Người ta làm một lối đi xung quanh vườn rộng $1\text{ m}$, phần đất còn lại bên trong dùng để trồng hoa.\par
1) Viết đa thức theo biến $x$ biểu thị diện tích mảnh vườn ban đầu.\par
2) Viết đa thức theo biến $x$ biểu thị diện tích phần đất bên trong dùng để trồng hoa.\par
\vspace{0.4mm}
\noindent\textit{Bài giải:}\dongke{9}
\vspace{1mm}
\textbf{b) (1,0 điểm)} Cho hai số $x, y$ thỏa mãn $x + y = 6$ và $xy = 5$. Không tính cụ thể giá trị của $x$ và $y$, hãy tính giá trị của biểu thức $T = x^3 + y^3$.\par
\vspace{0.4mm}
\noindent\textit{Bài giải:}\dongke{7}
\\ \hline
\end{tabularx}


\newpage

% #############################################################################
% ========================= TOÁN 8 - TUẦN 03 - BUỔI 02: MÃ ĐỀ 104 ============
% ========================= MẶT TRƯỚC (TRANG 7) ===============================
% #############################################################################

\noindent
\begin{tabularx}{\linewidth}{|X|c|}
\hline
\textbf{PHIẾU KIỂM TRA ĐỊNH KỲ TOÁN 8 - TUẦN 03 - BUỔI 02} & \textbf{Mã đề: 104} \\
\textit{Chủ đề: Ôn tập Chương 1 - Đa thức} & \textit{Thời gian: 45 phút} \\
\hline
\end{tabularx}

\vspace{1mm}

\noindent
\begin{tabularx}{\linewidth}{|X|}
\hline
\rule[-2mm]{0pt}{7.5mm}\textbf{Họ và tên:} \textbf{\large Hồ Hoàng Long} \hfill \textbf{Mã HS:} \texttt{HOANGLONG2711} \quad \\
\hline
\end{tabularx}

\vspace{1mm}

\noindent
\begin{tabularx}{\linewidth}{|p{3.2cm}|X|}
\hline
\centering\textbf{ĐIỂM SỐ} & \centering\textbf{LỜI NHẬN XÉT CỦA THẦY CÔ} \tabularnewline
\hline
\rule{0pt}{1.3cm} & \tabularnewline
\hline
\end{tabularx}

\vspace{1.2mm}

\noindent\textbf{BẢNG ĐIỀN ĐÁP ÁN TRẮC NGHIỆM (Mỗi câu đúng 0,5 điểm):}
\vspace{0.8mm}

\noindent
\renewcommand{\arraystretch}{1.35}
\begin{tabularx}{\linewidth}{|c|*{5}{>{\centering\arraybackslash}X|}}
\hline
\textbf{Câu} & \textbf{1} & \textbf{2} & \textbf{3} & \textbf{4} & \textbf{5} \\
\hline
\textbf{Đáp án} & & & & & \\
\hline
\end{tabularx}
\renewcommand{\arraystretch}{1.0}

\vspace{1.5mm}

\noindent\textbf{\large PHẦN I. TRẮC NGHIỆM KHÁCH QUAN (2,5 điểm)}
\vspace{1mm}

\cauTN{1}{Bậc của đa thức $Q = 5x^4 y^3 - 2x^3 y^4 + 4x^6 y - 5x^4 y^3 + 7$ là:}{$6$}{$5$}{$7$}{$4$}

\cauTN{2}{Đơn thức nào sau đây đồng dạng với đơn thức $-7x^4 y^2$?}{$3x^4 y^2$}{$-7x^2 y^4$}{$4x^4 y^3$}{$2x^3 y^2$}

\cauTN{3}{Kết quả của phép nhân đa thức $(3x + 2)(x - 2)$ là:}{$3x^2 + 4x - 4$}{$3x^2 - 4x - 4$}{$3x^2 - 6$}{$3x^2 - 8x - 4$}

\cauTN{4}{Kết quả của phép chia $(24x^5 y^3 - 18x^3 y^2) : 6x^3 y^2$ là:}{$4x^2 y - 3$}{$4x^2 y - 3xy$}{$4x^2 y^2 - 3$}{$4x^2 - 3y$}

\cauTN{5}{Giá trị của biểu thức $N = x^2 + 8xy + 16y^2$ tại $x = 12$ và $y = 2$ là:}{$256$}{$400$}{$324$}{$100$}

\vspace{1mm}

\noindent\textbf{\large PHẦN II. TỰ LUẬN (7,5 điểm)}

\noindent\textbf{Câu 1 (2,5 điểm). Thực hiện phép tính với đa thức:}
\vspace{0.8mm}

\noindent
\begin{tabularx}{\linewidth}{|X|X|}
\hline
\multicolumn{2}{|p{0.97\linewidth}|}{Cho hai đa thức: $A = 5x^2 y - 4xy^2 + 3xy - 1$ \ và \ $B = 3x^2 y + 4xy^2 - 2xy + 6$.} \\
\hline
\oBaiTap[7]{a)}{Tính $C = A + B$:} & 
\oBaiTap[7]{b)}{Tính $D = A - B$:} \\
\hline
\oBaiTap[5]{c)}{Tìm bậc của đa thức $C$:} & 
\oBaiTap[5]{d)}{Tính giá trị của đa thức $C$ tại $x = 1; y = -1$:} \\
\hline
\end{tabularx}


\newpage

% #############################################################################
% ========================= TOÁN 8 - TUẦN 03 - BUỔI 02: MÃ ĐỀ 104 ============
% ========================= MẶT SAU (TRANG 8) =================================
% #############################################################################

\noindent
\begin{tabularx}{\linewidth}{|X|c|}
\hline
\textbf{PHẦN II. TỰ LUẬN (Tiếp theo) - TOÁN 8 TUẦN 03 BUỔI 02} & \textbf{Mã đề: 104} \\
\hline
\end{tabularx}

\vspace{1mm}

\noindent
\begin{tabularx}{\linewidth}{|X|}
\hline
\rule[-2mm]{0pt}{7mm}\textbf{Họ và tên:} \textbf{\large Hồ Hoàng Long} \hfill \textbf{Mã HS:} \texttt{HOANGLONG2711} \quad \\
\hline
\end{tabularx}

\vspace{1.5mm}

\noindent\textbf{\large Câu 2 (2,5 điểm). Rút gọn và tìm $x$:}
\vspace{0.8mm}

\noindent
\begin{tabularx}{\linewidth}{|X|X|}
\hline
\oBaiTap[12]{a)}{Rút gọn và chứng minh giá trị biểu thức không phụ thuộc vào giá trị của biến $x$:\par
$N = (4x - 1)(x + 2) - 4x(x + 1) - 3x + 9$.} & 
\oBaiTap[12]{b)}{Tìm số thực $x$, biết:\par
$(2x - 1)(4x + 3) - 8x(x - 1) = 29$.} \\
\hline
\end{tabularx}

\vspace{1.8mm}

\noindent\textbf{\large Câu 3 (2,5 điểm). Bài toán thực tế và nâng cao:}
\vspace{0.8mm}

\noindent
\begin{tabularx}{\linewidth}{|X|}
\hline
\vspace{0.5mm}
\textbf{a) (1,5 điểm)} Một khu đất hình chữ nhật có chiều dài là $(4x + 5)\text{ m}$ và chiều rộng là $(2x + 4)\text{ m}$. Người ta làm một lối đi xung quanh khu đất rộng $1\text{ m}$, phần đất còn lại bên trong dùng để trồng rau sạch.\par
1) Viết đa thức theo biến $x$ biểu thị diện tích khu đất ban đầu.\par
2) Viết đa thức theo biến $x$ biểu thị diện tích phần đất bên trong dùng để trồng rau.\par
\vspace{0.4mm}
\noindent\textit{Bài giải:}\dongke{9}
\vspace{1mm}
\textbf{b) (1,0 điểm)} Cho hai số $x, y$ thỏa mãn $x + y = 6$ và $xy = 8$. Không tính cụ thể giá trị của $x$ và $y$, hãy tính giá trị của biểu thức $M = x^3 + y^3$.\par
\vspace{0.4mm}
\noindent\textit{Bài giải:}\dongke{7}
\\ \hline
\end{tabularx}



\newpage


% #############################################################################
% ========================= TOÁN 8 - TUẦN 03 - BUỔI 02: MÃ ĐỀ 101 ============
% ========================= MẶT TRƯỚC (TRANG 1) ===============================
% #############################################################################

\noindent
\begin{tabularx}{\linewidth}{|X|c|}
\hline
\textbf{PHIẾU KIỂM TRA ĐỊNH KỲ TOÁN 8 - TUẦN 03 - BUỔI 02} & \textbf{Mã đề: 101} \\
\textit{Chủ đề: Ôn tập Chương 1 - Đa thức} & \textit{Thời gian: 45 phút} \\
\hline
\end{tabularx}

\vspace{1mm}

\noindent
\begin{tabularx}{\linewidth}{|X|}
\hline
\rule[-2mm]{0pt}{7.5mm}\textbf{Họ và tên:} \textbf{\large Lê Nhật Minh} \hfill \textbf{Mã HS:} \texttt{NHATMINH1711} \quad \\
\hline
\end{tabularx}

\vspace{1mm}

\noindent
\begin{tabularx}{\linewidth}{|p{3.2cm}|X|}
\hline
\centering\textbf{ĐIỂM SỐ} & \centering\textbf{LỜI NHẬN XÉT CỦA THẦY CÔ} \tabularnewline
\hline
\rule{0pt}{1.3cm} & \tabularnewline
\hline
\end{tabularx}

\vspace{1.2mm}

\noindent\textbf{BẢNG ĐIỀN ĐÁP ÁN TRẮC NGHIỆM (Mỗi câu đúng 0,5 điểm):}
\vspace{0.8mm}

\noindent
\renewcommand{\arraystretch}{1.35}
\begin{tabularx}{\linewidth}{|c|*{5}{>{\centering\arraybackslash}X|}}
\hline
\textbf{Câu} & \textbf{1} & \textbf{2} & \textbf{3} & \textbf{4} & \textbf{5} \\
\hline
\textbf{Đáp án} & & & & & \\
\hline
\end{tabularx}
\renewcommand{\arraystretch}{1.0}

\vspace{1.5mm}

\noindent\textbf{\large PHẦN I. TRẮC NGHIỆM KHÁCH QUAN (2,5 điểm)}
\vspace{1mm}

\cauTN{1}{Bậc của đa thức $M = 5x^3 y^2 - 2x^2 y^3 + 3x^4 y - 5x^3 y^2 + 1$ là:}{$4$}{$5$}{$6$}{$3$}

\cauTN{2}{Đơn thức nào sau đây đồng dạng với đơn thức $-4x^3 y^2$?}{$2x^2 y^3$}{$-4x^3 y$}{$7x^3 y^2$}{$5x^2 y^2$}

\cauTN{3}{Kết quả của phép nhân đa thức $(2x - 3)(x + 2)$ là:}{$2x^2 + x - 6$}{$2x^2 - x - 6$}{$2x^2 + 7x - 6$}{$2x^2 - 6$}

\cauTN{4}{Kết quả của phép chia $(18x^4 y^3 - 12x^3 y^2) : 6x^2 y^2$ là:}{$3x^2 y - 2xy$}{$3x^2 y - 2$}{$3x^2 y^2 - 2x$}{$3x^2 y - 2x$}

\cauTN{5}{Giá trị của biểu thức $P = x^2 - 6xy + 9y^2$ tại $x = 16$ và $y = 2$ là:}{$100$}{$144$}{$64$}{$400$}

\vspace{1mm}

\noindent\textbf{\large PHẦN II. TỰ LUẬN (7,5 điểm)}

\noindent\textbf{Câu 1 (2,5 điểm). Thực hiện phép tính với đa thức:}
\vspace{0.8mm}

\noindent
\begin{tabularx}{\linewidth}{|X|X|}
\hline
\multicolumn{2}{|p{0.97\linewidth}|}{Cho hai đa thức: $A = 3x^2 y - 2xy^2 + 5xy - 1$ \ và \ $B = 2x^2 y + 2xy^2 - 3xy + 4$.} \\
\hline
\oBaiTap[7]{a)}{Tính $C = A + B$:} & 
\oBaiTap[7]{b)}{Tính $D = A - B$:} \\
\hline
\oBaiTap[5]{c)}{Tìm bậc của đa thức $C$:} & 
\oBaiTap[5]{d)}{Tính giá trị của đa thức $C$ tại $x = 1; y = -1$:} \\
\hline
\end{tabularx}


\newpage

% #############################################################################
% ========================= TOÁN 8 - TUẦN 03 - BUỔI 02: MÃ ĐỀ 101 ============
% ========================= MẶT SAU (TRANG 2) =================================
% #############################################################################

\noindent
\begin{tabularx}{\linewidth}{|X|c|}
\hline
\textbf{PHẦN II. TỰ LUẬN (Tiếp theo) - TOÁN 8 TUẦN 03 BUỔI 02} & \textbf{Mã đề: 101} \\
\hline
\end{tabularx}

\vspace{1mm}

\noindent
\begin{tabularx}{\linewidth}{|X|}
\hline
\rule[-2mm]{0pt}{7mm}\textbf{Họ và tên:} \textbf{\large Lê Nhật Minh} \hfill \textbf{Mã HS:} \texttt{NHATMINH1711} \quad \\
\hline
\end{tabularx}

\vspace{1.5mm}

\noindent\textbf{\large Câu 2 (2,5 điểm). Rút gọn và tìm $x$:}
\vspace{0.8mm}

\noindent
\begin{tabularx}{\linewidth}{|X|X|}
\hline
\oBaiTap[12]{a)}{Rút gọn và chứng minh giá trị biểu thức không phụ thuộc vào giá trị của biến $x$:\par
$M = (2x - 1)(x + 3) - 2x(x + 2) - x + 7$.} & 
\oBaiTap[12]{b)}{Tìm số thực $x$, biết:\par
$(3x - 1)(2x + 1) - 6x(x - 2) = 25$.} \\
\hline
\end{tabularx}

\vspace{1.8mm}

\noindent\textbf{\large Câu 3 (2,5 điểm). Bài toán thực tế và nâng cao:}
\vspace{0.8mm}

\noindent
\begin{tabularx}{\linewidth}{|X|}
\hline
\vspace{0.5mm}
\textbf{a) (1,5 điểm)} Một mảnh vườn hình chữ nhật có chiều dài là $(3x + 4)\text{ m}$ và chiều rộng là $(2x + 1)\text{ m}$. Người ta làm một lối đi xung quanh vườn rộng $1\text{ m}$, phần đất còn lại bên trong dùng để trồng hoa.\par
1) Viết đa thức theo biến $x$ biểu thị diện tích mảnh vườn ban đầu.\par
2) Viết đa thức theo biến $x$ biểu thị diện tích phần đất bên trong dùng để trồng hoa.\par
\vspace{0.4mm}
\noindent\textit{Bài giải:}\dongke{9}
\vspace{1mm}
\textbf{b) (1,0 điểm)} Cho hai số $x, y$ thỏa mãn $x + y = 4$ và $xy = 3$. Không tính cụ thể giá trị của $x$ và $y$, hãy tính giá trị của biểu thức $T = x^3 + y^3$.\par
\vspace{0.4mm}
\noindent\textit{Bài giải:}\dongke{7}
\\ \hline
\end{tabularx}


\newpage

% #############################################################################
% ========================= TOÁN 8 - TUẦN 03 - BUỔI 02: MÃ ĐỀ 102 ============
% ========================= MẶT TRƯỚC (TRANG 3) ===============================
% #############################################################################

\noindent
\begin{tabularx}{\linewidth}{|X|c|}
\hline
\textbf{PHIẾU KIỂM TRA ĐỊNH KỲ TOÁN 8 - TUẦN 03 - BUỔI 02} & \textbf{Mã đề: 102} \\
\textit{Chủ đề: Ôn tập Chương 1 - Đa thức} & \textit{Thời gian: 45 phút} \\
\hline
\end{tabularx}

\vspace{1mm}

\noindent
\begin{tabularx}{\linewidth}{|X|}
\hline
\rule[-2mm]{0pt}{7.5mm}\textbf{Họ và tên:} \textbf{\large Vũ Quang Minh} \hfill \textbf{Mã HS:} \texttt{QUANGMINH2601} \quad \\
\hline
\end{tabularx}

\vspace{1mm}

\noindent
\begin{tabularx}{\linewidth}{|p{3.2cm}|X|}
\hline
\centering\textbf{ĐIỂM SỐ} & \centering\textbf{LỜI NHẬN XÉT CỦA THẦY CÔ} \tabularnewline
\hline
\rule{0pt}{1.3cm} & \tabularnewline
\hline
\end{tabularx}

\vspace{1.2mm}

\noindent\textbf{BẢNG ĐIỀN ĐÁP ÁN TRẮC NGHIỆM (Mỗi câu đúng 0,5 điểm):}
\vspace{0.8mm}

\noindent
\renewcommand{\arraystretch}{1.35}
\begin{tabularx}{\linewidth}{|c|*{5}{>{\centering\arraybackslash}X|}}
\hline
\textbf{Câu} & \textbf{1} & \textbf{2} & \textbf{3} & \textbf{4} & \textbf{5} \\
\hline
\textbf{Đáp án} & & & & & \\
\hline
\end{tabularx}
\renewcommand{\arraystretch}{1.0}

\vspace{1.5mm}

\noindent\textbf{\large PHẦN I. TRẮC NGHIỆM KHÁCH QUAN (2,5 điểm)}
\vspace{1mm}

\cauTN{1}{Bậc của đa thức $N = 4x^3 y^3 - 3x^2 y^4 + 2x^5 y - 4x^3 y^3 + 9$ là:}{$5$}{$7$}{$6$}{$4$}

\cauTN{2}{Đơn thức nào sau đây đồng dạng với đơn thức $-5x^2 y^4$?}{$3x^2 y^4$}{$-5x^4 y^2$}{$2x^2 y^2$}{$4xy^4$}

\cauTN{3}{Kết quả của phép nhân đa thức $(2x + 3)(x - 2)$ là:}{$2x^2 + x - 6$}{$2x^2 - 7x - 6$}{$2x^2 - 6$}{$2x^2 - x - 6$}

\cauTN{4}{Kết quả của phép chia $(20x^5 y^3 - 15x^3 y^2) : 5x^3 y^2$ là:}{$4x^2 y - 3$}{$4x^2 y - 3xy$}{$4x^2 y^2 - 3$}{$4x^2 - 3y$}

\cauTN{5}{Giá trị của biểu thức $Q = x^2 + 6xy + 9y^2$ tại $x = 14$ và $y = 2$ là:}{$256$}{$400$}{$324$}{$100$}

\vspace{1mm}

\noindent\textbf{\large PHẦN II. TỰ LUẬN (7,5 điểm)}

\noindent\textbf{Câu 1 (2,5 điểm). Thực hiện phép tính với đa thức:}
\vspace{0.8mm}

\noindent
\begin{tabularx}{\linewidth}{|X|X|}
\hline
\multicolumn{2}{|p{0.97\linewidth}|}{Cho hai đa thức: $A = 4x^2 y - 3xy^2 + 4xy - 2$ \ và \ $B = 3x^2 y + 3xy^2 - 2xy + 5$.} \\
\hline
\oBaiTap[7]{a)}{Tính $C = A + B$:} & 
\oBaiTap[7]{b)}{Tính $D = A - B$:} \\
\hline
\oBaiTap[5]{c)}{Tìm bậc của đa thức $C$:} & 
\oBaiTap[5]{d)}{Tính giá trị của đa thức $C$ tại $x = 1; y = -1$:} \\
\hline
\end{tabularx}


\newpage

% #############################################################################
% ========================= TOÁN 8 - TUẦN 03 - BUỔI 02: MÃ ĐỀ 102 ============
% ========================= MẶT SAU (TRANG 4) =================================
% #############################################################################

\noindent
\begin{tabularx}{\linewidth}{|X|c|}
\hline
\textbf{PHẦN II. TỰ LUẬN (Tiếp theo) - TOÁN 8 TUẦN 03 BUỔI 02} & \textbf{Mã đề: 102} \\
\hline
\end{tabularx}

\vspace{1mm}

\noindent
\begin{tabularx}{\linewidth}{|X|}
\hline
\rule[-2mm]{0pt}{7mm}\textbf{Họ và tên:} \textbf{\large Vũ Quang Minh} \hfill \textbf{Mã HS:} \texttt{QUANGMINH2601} \quad \\
\hline
\end{tabularx}

\vspace{1.5mm}

\noindent\textbf{\large Câu 2 (2,5 điểm). Rút gọn và tìm $x$:}
\vspace{0.8mm}

\noindent
\begin{tabularx}{\linewidth}{|X|X|}
\hline
\oBaiTap[12]{a)}{Rút gọn và chứng minh giá trị biểu thức không phụ thuộc vào giá trị của biến $x$:\par
$N = (3x - 1)(x + 2) - 3x(x + 1) - 2x + 5$.} & 
\oBaiTap[12]{b)}{Tìm số thực $x$, biết:\par
$(2x - 1)(3x + 2) - 6x(x - 1) = 19$.} \\
\hline
\end{tabularx}

\vspace{1.8mm}

\noindent\textbf{\large Câu 3 (2,5 điểm). Bài toán thực tế và nâng cao:}
\vspace{0.8mm}

\noindent
\begin{tabularx}{\linewidth}{|X|}
\hline
\vspace{0.5mm}
\textbf{a) (1,5 điểm)} Một khu đất hình chữ nhật có chiều dài là $(4x + 3)\text{ m}$ và chiều rộng là $(2x + 3)\text{ m}$. Người ta làm một lối đi xung quanh khu đất rộng $1\text{ m}$, phần đất còn lại bên trong dùng để trồng rau sạch.\par
1) Viết đa thức theo biến $x$ biểu thị diện tích khu đất ban đầu.\par
2) Viết đa thức theo biến $x$ biểu thị diện tích phần đất bên trong dùng để trồng rau.\par
\vspace{0.4mm}
\noindent\textit{Bài giải:}\dongke{9}
\vspace{1mm}
\textbf{b) (1,0 điểm)} Cho hai số $x, y$ thỏa mãn $x + y = 5$ và $xy = 4$. Không tính cụ thể giá trị của $x$ và $y$, hãy tính giá trị của biểu thức $M = x^3 + y^3$.\par
\vspace{0.4mm}
\noindent\textit{Bài giải:}\dongke{7}
\\ \hline
\end{tabularx}


\newpage

% #############################################################################
% ========================= TOÁN 8 - TUẦN 03 - BUỔI 02: MÃ ĐỀ 103 ============
% ========================= MẶT TRƯỚC (TRANG 5) ===============================
% #############################################################################

\noindent
\begin{tabularx}{\linewidth}{|X|c|}
\hline
\textbf{PHIẾU KIỂM TRA ĐỊNH KỲ TOÁN 8 - TUẦN 03 - BUỔI 02} & \textbf{Mã đề: 103} \\
\textit{Chủ đề: Ôn tập Chương 1 - Đa thức} & \textit{Thời gian: 45 phút} \\
\hline
\end{tabularx}

\vspace{1mm}

\noindent
\begin{tabularx}{\linewidth}{|X|}
\hline
\rule[-2mm]{0pt}{7.5mm}\textbf{Họ và tên:} \textbf{\large Đinh Mộc Miên} \hfill \textbf{Mã HS:} \texttt{MOCMIEN1111} \quad \\
\hline
\end{tabularx}

\vspace{1mm}

\noindent
\begin{tabularx}{\linewidth}{|p{3.2cm}|X|}
\hline
\centering\textbf{ĐIỂM SỐ} & \centering\textbf{LỜI NHẬN XÉT CỦA THẦY CÔ} \tabularnewline
\hline
\rule{0pt}{1.3cm} & \tabularnewline
\hline
\end{tabularx}

\vspace{1.2mm}

\noindent\textbf{BẢNG ĐIỀN ĐÁP ÁN TRẮC NGHIỆM (Mỗi câu đúng 0,5 điểm):}
\vspace{0.8mm}

\noindent
\renewcommand{\arraystretch}{1.35}
\begin{tabularx}{\linewidth}{|c|*{5}{>{\centering\arraybackslash}X|}}
\hline
\textbf{Câu} & \textbf{1} & \textbf{2} & \textbf{3} & \textbf{4} & \textbf{5} \\
\hline
\textbf{Đáp án} & & & & & \\
\hline
\end{tabularx}
\renewcommand{\arraystretch}{1.0}

\vspace{1.5mm}

\noindent\textbf{\large PHẦN I. TRẮC NGHIỆM KHÁCH QUAN (2,5 điểm)}
\vspace{1mm}

\cauTN{1}{Bậc của đa thức $P = 6x^4 y^2 - 3x^3 y^3 + 2x^5 y - 6x^4 y^2 + 8$ là:}{$5$}{$6$}{$7$}{$4$}

\cauTN{2}{Đơn thức nào sau đây đồng dạng với đơn thức $-6x^3 y^3$?}{$5x^3 y^3$}{$-6x^2 y^3$}{$3x^3 y^2$}{$4x^4 y^2$}

\cauTN{3}{Kết quả của phép nhân đa thức $(3x - 2)(x + 2)$ là:}{$3x^2 + 4x - 4$}{$3x^2 - 4x - 4$}{$3x^2 + 6x - 4$}{$3x^2 - 4$}

\cauTN{4}{Kết quả của phép chia $(24x^4 y^3 - 16x^3 y^2) : 8x^2 y^2$ là:}{$3x^2 y - 2x$}{$3x^2 y - 2$}{$3x^2 y^2 - 2x$}{$3x^2 - 2y$}

\cauTN{5}{Giá trị của biểu thức $M = x^2 - 8xy + 16y^2$ tại $x = 18$ và $y = 2$ là:}{$64$}{$144$}{$100$}{$400$}

\vspace{1mm}

\noindent\textbf{\large PHẦN II. TỰ LUẬN (7,5 điểm)}

\noindent\textbf{Câu 1 (2,5 điểm). Thực hiện phép tính với đa thức:}
\vspace{0.8mm}

\noindent
\begin{tabularx}{\linewidth}{|X|X|}
\hline
\multicolumn{2}{|p{0.97\linewidth}|}{Cho hai đa thức: $A = 5x^2 y - 3xy^2 + 6xy - 3$ \ và \ $B = 2x^2 y + 3xy^2 - 4xy + 5$.} \\
\hline
\oBaiTap[7]{a)}{Tính $C = A + B$:} & 
\oBaiTap[7]{b)}{Tính $D = A - B$:} \\
\hline
\oBaiTap[5]{c)}{Tìm bậc của đa thức $C$:} & 
\oBaiTap[5]{d)}{Tính giá trị của đa thức $C$ tại $x = 1; y = -1$:} \\
\hline
\end{tabularx}


\newpage

% #############################################################################
% ========================= TOÁN 8 - TUẦN 03 - BUỔI 02: MÃ ĐỀ 103 ============
% ========================= MẶT SAU (TRANG 6) =================================
% #############################################################################

\noindent
\begin{tabularx}{\linewidth}{|X|c|}
\hline
\textbf{PHẦN II. TỰ LUẬN (Tiếp theo) - TOÁN 8 TUẦN 03 BUỔI 02} & \textbf{Mã đề: 103} \\
\hline
\end{tabularx}

\vspace{1mm}

\noindent
\begin{tabularx}{\linewidth}{|X|}
\hline
\rule[-2mm]{0pt}{7mm}\textbf{Họ và tên:} \textbf{\large Đinh Mộc Miên} \hfill \textbf{Mã HS:} \texttt{MOCMIEN1111} \quad \\
\hline
\end{tabularx}

\vspace{1.5mm}

\noindent\textbf{\large Câu 2 (2,5 điểm). Rút gọn và tìm $x$:}
\vspace{0.8mm}

\noindent
\begin{tabularx}{\linewidth}{|X|X|}
\hline
\oBaiTap[12]{a)}{Rút gọn và chứng minh giá trị biểu thức không phụ thuộc vào giá trị của biến $x$:\par
$P = (2x - 3)(x + 2) - 2x(x + 1) + x + 8$.} & 
\oBaiTap[12]{b)}{Tìm số thực $x$, biết:\par
$(3x - 2)(2x + 1) - 6x(x - 1) = 23$.} \\
\hline
\end{tabularx}

\vspace{1.8mm}

\noindent\textbf{\large Câu 3 (2,5 điểm). Bài toán thực tế và nâng cao:}
\vspace{0.8mm}

\noindent
\begin{tabularx}{\linewidth}{|X|}
\hline
\vspace{0.5mm}
\textbf{a) (1,5 điểm)} Một mảnh vườn hình chữ nhật có chiều dài là $(3x + 5)\text{ m}$ và chiều rộng là $(2x + 2)\text{ m}$. Người ta làm một lối đi xung quanh vườn rộng $1\text{ m}$, phần đất còn lại bên trong dùng để trồng hoa.\par
1) Viết đa thức theo biến $x$ biểu thị diện tích mảnh vườn ban đầu.\par
2) Viết đa thức theo biến $x$ biểu thị diện tích phần đất bên trong dùng để trồng hoa.\par
\vspace{0.4mm}
\noindent\textit{Bài giải:}\dongke{9}
\vspace{1mm}
\textbf{b) (1,0 điểm)} Cho hai số $x, y$ thỏa mãn $x + y = 6$ và $xy = 5$. Không tính cụ thể giá trị của $x$ và $y$, hãy tính giá trị của biểu thức $T = x^3 + y^3$.\par
\vspace{0.4mm}
\noindent\textit{Bài giải:}\dongke{7}
\\ \hline
\end{tabularx}


\newpage

% #############################################################################
% ========================= TOÁN 8 - TUẦN 03 - BUỔI 02: MÃ ĐỀ 104 ============
% ========================= MẶT TRƯỚC (TRANG 7) ===============================
% #############################################################################

\noindent
\begin{tabularx}{\linewidth}{|X|c|}
\hline
\textbf{PHIẾU KIỂM TRA ĐỊNH KỲ TOÁN 8 - TUẦN 03 - BUỔI 02} & \textbf{Mã đề: 104} \\
\textit{Chủ đề: Ôn tập Chương 1 - Đa thức} & \textit{Thời gian: 45 phút} \\
\hline
\end{tabularx}

\vspace{1mm}

\noindent
\begin{tabularx}{\linewidth}{|X|}
\hline
\rule[-2mm]{0pt}{7.5mm}\textbf{Họ và tên:} \textbf{\large Nguyễn Đình Nam} \hfill \textbf{Mã HS:} \texttt{DINHNAM2510} \quad \\
\hline
\end{tabularx}

\vspace{1mm}

\noindent
\begin{tabularx}{\linewidth}{|p{3.2cm}|X|}
\hline
\centering\textbf{ĐIỂM SỐ} & \centering\textbf{LỜI NHẬN XÉT CỦA THẦY CÔ} \tabularnewline
\hline
\rule{0pt}{1.3cm} & \tabularnewline
\hline
\end{tabularx}

\vspace{1.2mm}

\noindent\textbf{BẢNG ĐIỀN ĐÁP ÁN TRẮC NGHIỆM (Mỗi câu đúng 0,5 điểm):}
\vspace{0.8mm}

\noindent
\renewcommand{\arraystretch}{1.35}
\begin{tabularx}{\linewidth}{|c|*{5}{>{\centering\arraybackslash}X|}}
\hline
\textbf{Câu} & \textbf{1} & \textbf{2} & \textbf{3} & \textbf{4} & \textbf{5} \\
\hline
\textbf{Đáp án} & & & & & \\
\hline
\end{tabularx}
\renewcommand{\arraystretch}{1.0}

\vspace{1.5mm}

\noindent\textbf{\large PHẦN I. TRẮC NGHIỆM KHÁCH QUAN (2,5 điểm)}
\vspace{1mm}

\cauTN{1}{Bậc của đa thức $Q = 5x^4 y^3 - 2x^3 y^4 + 4x^6 y - 5x^4 y^3 + 7$ là:}{$6$}{$5$}{$7$}{$4$}

\cauTN{2}{Đơn thức nào sau đây đồng dạng với đơn thức $-7x^4 y^2$?}{$3x^4 y^2$}{$-7x^2 y^4$}{$4x^4 y^3$}{$2x^3 y^2$}

\cauTN{3}{Kết quả của phép nhân đa thức $(3x + 2)(x - 2)$ là:}{$3x^2 + 4x - 4$}{$3x^2 - 4x - 4$}{$3x^2 - 6$}{$3x^2 - 8x - 4$}

\cauTN{4}{Kết quả của phép chia $(24x^5 y^3 - 18x^3 y^2) : 6x^3 y^2$ là:}{$4x^2 y - 3$}{$4x^2 y - 3xy$}{$4x^2 y^2 - 3$}{$4x^2 - 3y$}

\cauTN{5}{Giá trị của biểu thức $N = x^2 + 8xy + 16y^2$ tại $x = 12$ và $y = 2$ là:}{$256$}{$400$}{$324$}{$100$}

\vspace{1mm}

\noindent\textbf{\large PHẦN II. TỰ LUẬN (7,5 điểm)}

\noindent\textbf{Câu 1 (2,5 điểm). Thực hiện phép tính với đa thức:}
\vspace{0.8mm}

\noindent
\begin{tabularx}{\linewidth}{|X|X|}
\hline
\multicolumn{2}{|p{0.97\linewidth}|}{Cho hai đa thức: $A = 5x^2 y - 4xy^2 + 3xy - 1$ \ và \ $B = 3x^2 y + 4xy^2 - 2xy + 6$.} \\
\hline
\oBaiTap[7]{a)}{Tính $C = A + B$:} & 
\oBaiTap[7]{b)}{Tính $D = A - B$:} \\
\hline
\oBaiTap[5]{c)}{Tìm bậc của đa thức $C$:} & 
\oBaiTap[5]{d)}{Tính giá trị của đa thức $C$ tại $x = 1; y = -1$:} \\
\hline
\end{tabularx}


\newpage

% #############################################################################
% ========================= TOÁN 8 - TUẦN 03 - BUỔI 02: MÃ ĐỀ 104 ============
% ========================= MẶT SAU (TRANG 8) =================================
% #############################################################################

\noindent
\begin{tabularx}{\linewidth}{|X|c|}
\hline
\textbf{PHẦN II. TỰ LUẬN (Tiếp theo) - TOÁN 8 TUẦN 03 BUỔI 02} & \textbf{Mã đề: 104} \\
\hline
\end{tabularx}

\vspace{1mm}

\noindent
\begin{tabularx}{\linewidth}{|X|}
\hline
\rule[-2mm]{0pt}{7mm}\textbf{Họ và tên:} \textbf{\large Nguyễn Đình Nam} \hfill \textbf{Mã HS:} \texttt{DINHNAM2510} \quad \\
\hline
\end{tabularx}

\vspace{1.5mm}

\noindent\textbf{\large Câu 2 (2,5 điểm). Rút gọn và tìm $x$:}
\vspace{0.8mm}

\noindent
\begin{tabularx}{\linewidth}{|X|X|}
\hline
\oBaiTap[12]{a)}{Rút gọn và chứng minh giá trị biểu thức không phụ thuộc vào giá trị của biến $x$:\par
$N = (4x - 1)(x + 2) - 4x(x + 1) - 3x + 9$.} & 
\oBaiTap[12]{b)}{Tìm số thực $x$, biết:\par
$(2x - 1)(4x + 3) - 8x(x - 1) = 29$.} \\
\hline
\end{tabularx}

\vspace{1.8mm}

\noindent\textbf{\large Câu 3 (2,5 điểm). Bài toán thực tế và nâng cao:}
\vspace{0.8mm}

\noindent
\begin{tabularx}{\linewidth}{|X|}
\hline
\vspace{0.5mm}
\textbf{a) (1,5 điểm)} Một khu đất hình chữ nhật có chiều dài là $(4x + 5)\text{ m}$ và chiều rộng là $(2x + 4)\text{ m}$. Người ta làm một lối đi xung quanh khu đất rộng $1\text{ m}$, phần đất còn lại bên trong dùng để trồng rau sạch.\par
1) Viết đa thức theo biến $x$ biểu thị diện tích khu đất ban đầu.\par
2) Viết đa thức theo biến $x$ biểu thị diện tích phần đất bên trong dùng để trồng rau.\par
\vspace{0.4mm}
\noindent\textit{Bài giải:}\dongke{9}
\vspace{1mm}
\textbf{b) (1,0 điểm)} Cho hai số $x, y$ thỏa mãn $x + y = 6$ và $xy = 8$. Không tính cụ thể giá trị của $x$ và $y$, hãy tính giá trị của biểu thức $M = x^3 + y^3$.\par
\vspace{0.4mm}
\noindent\textit{Bài giải:}\dongke{7}
\\ \hline
\end{tabularx}



\newpage


% #############################################################################
% ========================= TOÁN 8 - TUẦN 03 - BUỔI 02: MÃ ĐỀ 101 ============
% ========================= MẶT TRƯỚC (TRANG 1) ===============================
% #############################################################################

\noindent
\begin{tabularx}{\linewidth}{|X|c|}
\hline
\textbf{PHIẾU KIỂM TRA ĐỊNH KỲ TOÁN 8 - TUẦN 03 - BUỔI 02} & \textbf{Mã đề: 101} \\
\textit{Chủ đề: Ôn tập Chương 1 - Đa thức} & \textit{Thời gian: 45 phút} \\
\hline
\end{tabularx}

\vspace{1mm}

\noindent
\begin{tabularx}{\linewidth}{|X|}
\hline
\rule[-2mm]{0pt}{7.5mm}\textbf{Họ và tên:} \textbf{\large Lê Thái Nguyên} \hfill \textbf{Mã HS:} \texttt{THAINGUYEN1112} \quad \\
\hline
\end{tabularx}

\vspace{1mm}

\noindent
\begin{tabularx}{\linewidth}{|p{3.2cm}|X|}
\hline
\centering\textbf{ĐIỂM SỐ} & \centering\textbf{LỜI NHẬN XÉT CỦA THẦY CÔ} \tabularnewline
\hline
\rule{0pt}{1.3cm} & \tabularnewline
\hline
\end{tabularx}

\vspace{1.2mm}

\noindent\textbf{BẢNG ĐIỀN ĐÁP ÁN TRẮC NGHIỆM (Mỗi câu đúng 0,5 điểm):}
\vspace{0.8mm}

\noindent
\renewcommand{\arraystretch}{1.35}
\begin{tabularx}{\linewidth}{|c|*{5}{>{\centering\arraybackslash}X|}}
\hline
\textbf{Câu} & \textbf{1} & \textbf{2} & \textbf{3} & \textbf{4} & \textbf{5} \\
\hline
\textbf{Đáp án} & & & & & \\
\hline
\end{tabularx}
\renewcommand{\arraystretch}{1.0}

\vspace{1.5mm}

\noindent\textbf{\large PHẦN I. TRẮC NGHIỆM KHÁCH QUAN (2,5 điểm)}
\vspace{1mm}

\cauTN{1}{Bậc của đa thức $M = 5x^3 y^2 - 2x^2 y^3 + 3x^4 y - 5x^3 y^2 + 1$ là:}{$4$}{$5$}{$6$}{$3$}

\cauTN{2}{Đơn thức nào sau đây đồng dạng với đơn thức $-4x^3 y^2$?}{$2x^2 y^3$}{$-4x^3 y$}{$7x^3 y^2$}{$5x^2 y^2$}

\cauTN{3}{Kết quả của phép nhân đa thức $(2x - 3)(x + 2)$ là:}{$2x^2 + x - 6$}{$2x^2 - x - 6$}{$2x^2 + 7x - 6$}{$2x^2 - 6$}

\cauTN{4}{Kết quả của phép chia $(18x^4 y^3 - 12x^3 y^2) : 6x^2 y^2$ là:}{$3x^2 y - 2xy$}{$3x^2 y - 2$}{$3x^2 y^2 - 2x$}{$3x^2 y - 2x$}

\cauTN{5}{Giá trị của biểu thức $P = x^2 - 6xy + 9y^2$ tại $x = 16$ và $y = 2$ là:}{$100$}{$144$}{$64$}{$400$}

\vspace{1mm}

\noindent\textbf{\large PHẦN II. TỰ LUẬN (7,5 điểm)}

\noindent\textbf{Câu 1 (2,5 điểm). Thực hiện phép tính với đa thức:}
\vspace{0.8mm}

\noindent
\begin{tabularx}{\linewidth}{|X|X|}
\hline
\multicolumn{2}{|p{0.97\linewidth}|}{Cho hai đa thức: $A = 3x^2 y - 2xy^2 + 5xy - 1$ \ và \ $B = 2x^2 y + 2xy^2 - 3xy + 4$.} \\
\hline
\oBaiTap[7]{a)}{Tính $C = A + B$:} & 
\oBaiTap[7]{b)}{Tính $D = A - B$:} \\
\hline
\oBaiTap[5]{c)}{Tìm bậc của đa thức $C$:} & 
\oBaiTap[5]{d)}{Tính giá trị của đa thức $C$ tại $x = 1; y = -1$:} \\
\hline
\end{tabularx}


\newpage

% #############################################################################
% ========================= TOÁN 8 - TUẦN 03 - BUỔI 02: MÃ ĐỀ 101 ============
% ========================= MẶT SAU (TRANG 2) =================================
% #############################################################################

\noindent
\begin{tabularx}{\linewidth}{|X|c|}
\hline
\textbf{PHẦN II. TỰ LUẬN (Tiếp theo) - TOÁN 8 TUẦN 03 BUỔI 02} & \textbf{Mã đề: 101} \\
\hline
\end{tabularx}

\vspace{1mm}

\noindent
\begin{tabularx}{\linewidth}{|X|}
\hline
\rule[-2mm]{0pt}{7mm}\textbf{Họ và tên:} \textbf{\large Lê Thái Nguyên} \hfill \textbf{Mã HS:} \texttt{THAINGUYEN1112} \quad \\
\hline
\end{tabularx}

\vspace{1.5mm}

\noindent\textbf{\large Câu 2 (2,5 điểm). Rút gọn và tìm $x$:}
\vspace{0.8mm}

\noindent
\begin{tabularx}{\linewidth}{|X|X|}
\hline
\oBaiTap[12]{a)}{Rút gọn và chứng minh giá trị biểu thức không phụ thuộc vào giá trị của biến $x$:\par
$M = (2x - 1)(x + 3) - 2x(x + 2) - x + 7$.} & 
\oBaiTap[12]{b)}{Tìm số thực $x$, biết:\par
$(3x - 1)(2x + 1) - 6x(x - 2) = 25$.} \\
\hline
\end{tabularx}

\vspace{1.8mm}

\noindent\textbf{\large Câu 3 (2,5 điểm). Bài toán thực tế và nâng cao:}
\vspace{0.8mm}

\noindent
\begin{tabularx}{\linewidth}{|X|}
\hline
\vspace{0.5mm}
\textbf{a) (1,5 điểm)} Một mảnh vườn hình chữ nhật có chiều dài là $(3x + 4)\text{ m}$ và chiều rộng là $(2x + 1)\text{ m}$. Người ta làm một lối đi xung quanh vườn rộng $1\text{ m}$, phần đất còn lại bên trong dùng để trồng hoa.\par
1) Viết đa thức theo biến $x$ biểu thị diện tích mảnh vườn ban đầu.\par
2) Viết đa thức theo biến $x$ biểu thị diện tích phần đất bên trong dùng để trồng hoa.\par
\vspace{0.4mm}
\noindent\textit{Bài giải:}\dongke{9}
\vspace{1mm}
\textbf{b) (1,0 điểm)} Cho hai số $x, y$ thỏa mãn $x + y = 4$ và $xy = 3$. Không tính cụ thể giá trị của $x$ và $y$, hãy tính giá trị của biểu thức $T = x^3 + y^3$.\par
\vspace{0.4mm}
\noindent\textit{Bài giải:}\dongke{7}
\\ \hline
\end{tabularx}


\newpage

% #############################################################################
% ========================= TOÁN 8 - TUẦN 03 - BUỔI 02: MÃ ĐỀ 102 ============
% ========================= MẶT TRƯỚC (TRANG 3) ===============================
% #############################################################################

\noindent
\begin{tabularx}{\linewidth}{|X|c|}
\hline
\textbf{PHIẾU KIỂM TRA ĐỊNH KỲ TOÁN 8 - TUẦN 03 - BUỔI 02} & \textbf{Mã đề: 102} \\
\textit{Chủ đề: Ôn tập Chương 1 - Đa thức} & \textit{Thời gian: 45 phút} \\
\hline
\end{tabularx}

\vspace{1mm}

\noindent
\begin{tabularx}{\linewidth}{|X|}
\hline
\rule[-2mm]{0pt}{7.5mm}\textbf{Họ và tên:} \textbf{\large Lê Bảo Ngọc} \hfill \textbf{Mã HS:} \texttt{BAONGOC1402} \quad \\
\hline
\end{tabularx}

\vspace{1mm}

\noindent
\begin{tabularx}{\linewidth}{|p{3.2cm}|X|}
\hline
\centering\textbf{ĐIỂM SỐ} & \centering\textbf{LỜI NHẬN XÉT CỦA THẦY CÔ} \tabularnewline
\hline
\rule{0pt}{1.3cm} & \tabularnewline
\hline
\end{tabularx}

\vspace{1.2mm}

\noindent\textbf{BẢNG ĐIỀN ĐÁP ÁN TRẮC NGHIỆM (Mỗi câu đúng 0,5 điểm):}
\vspace{0.8mm}

\noindent
\renewcommand{\arraystretch}{1.35}
\begin{tabularx}{\linewidth}{|c|*{5}{>{\centering\arraybackslash}X|}}
\hline
\textbf{Câu} & \textbf{1} & \textbf{2} & \textbf{3} & \textbf{4} & \textbf{5} \\
\hline
\textbf{Đáp án} & & & & & \\
\hline
\end{tabularx}
\renewcommand{\arraystretch}{1.0}

\vspace{1.5mm}

\noindent\textbf{\large PHẦN I. TRẮC NGHIỆM KHÁCH QUAN (2,5 điểm)}
\vspace{1mm}

\cauTN{1}{Bậc của đa thức $N = 4x^3 y^3 - 3x^2 y^4 + 2x^5 y - 4x^3 y^3 + 9$ là:}{$5$}{$7$}{$6$}{$4$}

\cauTN{2}{Đơn thức nào sau đây đồng dạng với đơn thức $-5x^2 y^4$?}{$3x^2 y^4$}{$-5x^4 y^2$}{$2x^2 y^2$}{$4xy^4$}

\cauTN{3}{Kết quả của phép nhân đa thức $(2x + 3)(x - 2)$ là:}{$2x^2 + x - 6$}{$2x^2 - 7x - 6$}{$2x^2 - 6$}{$2x^2 - x - 6$}

\cauTN{4}{Kết quả của phép chia $(20x^5 y^3 - 15x^3 y^2) : 5x^3 y^2$ là:}{$4x^2 y - 3$}{$4x^2 y - 3xy$}{$4x^2 y^2 - 3$}{$4x^2 - 3y$}

\cauTN{5}{Giá trị của biểu thức $Q = x^2 + 6xy + 9y^2$ tại $x = 14$ và $y = 2$ là:}{$256$}{$400$}{$324$}{$100$}

\vspace{1mm}

\noindent\textbf{\large PHẦN II. TỰ LUẬN (7,5 điểm)}

\noindent\textbf{Câu 1 (2,5 điểm). Thực hiện phép tính với đa thức:}
\vspace{0.8mm}

\noindent
\begin{tabularx}{\linewidth}{|X|X|}
\hline
\multicolumn{2}{|p{0.97\linewidth}|}{Cho hai đa thức: $A = 4x^2 y - 3xy^2 + 4xy - 2$ \ và \ $B = 3x^2 y + 3xy^2 - 2xy + 5$.} \\
\hline
\oBaiTap[7]{a)}{Tính $C = A + B$:} & 
\oBaiTap[7]{b)}{Tính $D = A - B$:} \\
\hline
\oBaiTap[5]{c)}{Tìm bậc của đa thức $C$:} & 
\oBaiTap[5]{d)}{Tính giá trị của đa thức $C$ tại $x = 1; y = -1$:} \\
\hline
\end{tabularx}


\newpage

% #############################################################################
% ========================= TOÁN 8 - TUẦN 03 - BUỔI 02: MÃ ĐỀ 102 ============
% ========================= MẶT SAU (TRANG 4) =================================
% #############################################################################

\noindent
\begin{tabularx}{\linewidth}{|X|c|}
\hline
\textbf{PHẦN II. TỰ LUẬN (Tiếp theo) - TOÁN 8 TUẦN 03 BUỔI 02} & \textbf{Mã đề: 102} \\
\hline
\end{tabularx}

\vspace{1mm}

\noindent
\begin{tabularx}{\linewidth}{|X|}
\hline
\rule[-2mm]{0pt}{7mm}\textbf{Họ và tên:} \textbf{\large Lê Bảo Ngọc} \hfill \textbf{Mã HS:} \texttt{BAONGOC1402} \quad \\
\hline
\end{tabularx}

\vspace{1.5mm}

\noindent\textbf{\large Câu 2 (2,5 điểm). Rút gọn và tìm $x$:}
\vspace{0.8mm}

\noindent
\begin{tabularx}{\linewidth}{|X|X|}
\hline
\oBaiTap[12]{a)}{Rút gọn và chứng minh giá trị biểu thức không phụ thuộc vào giá trị của biến $x$:\par
$N = (3x - 1)(x + 2) - 3x(x + 1) - 2x + 5$.} & 
\oBaiTap[12]{b)}{Tìm số thực $x$, biết:\par
$(2x - 1)(3x + 2) - 6x(x - 1) = 19$.} \\
\hline
\end{tabularx}

\vspace{1.8mm}

\noindent\textbf{\large Câu 3 (2,5 điểm). Bài toán thực tế và nâng cao:}
\vspace{0.8mm}

\noindent
\begin{tabularx}{\linewidth}{|X|}
\hline
\vspace{0.5mm}
\textbf{a) (1,5 điểm)} Một khu đất hình chữ nhật có chiều dài là $(4x + 3)\text{ m}$ và chiều rộng là $(2x + 3)\text{ m}$. Người ta làm một lối đi xung quanh khu đất rộng $1\text{ m}$, phần đất còn lại bên trong dùng để trồng rau sạch.\par
1) Viết đa thức theo biến $x$ biểu thị diện tích khu đất ban đầu.\par
2) Viết đa thức theo biến $x$ biểu thị diện tích phần đất bên trong dùng để trồng rau.\par
\vspace{0.4mm}
\noindent\textit{Bài giải:}\dongke{9}
\vspace{1mm}
\textbf{b) (1,0 điểm)} Cho hai số $x, y$ thỏa mãn $x + y = 5$ và $xy = 4$. Không tính cụ thể giá trị của $x$ và $y$, hãy tính giá trị của biểu thức $M = x^3 + y^3$.\par
\vspace{0.4mm}
\noindent\textit{Bài giải:}\dongke{7}
\\ \hline
\end{tabularx}


\newpage

% #############################################################################
% ========================= TOÁN 8 - TUẦN 03 - BUỔI 02: MÃ ĐỀ 103 ============
% ========================= MẶT TRƯỚC (TRANG 5) ===============================
% #############################################################################

\noindent
\begin{tabularx}{\linewidth}{|X|c|}
\hline
\textbf{PHIẾU KIỂM TRA ĐỊNH KỲ TOÁN 8 - TUẦN 03 - BUỔI 02} & \textbf{Mã đề: 103} \\
\textit{Chủ đề: Ôn tập Chương 1 - Đa thức} & \textit{Thời gian: 45 phút} \\
\hline
\end{tabularx}

\vspace{1mm}

\noindent
\begin{tabularx}{\linewidth}{|X|}
\hline
\rule[-2mm]{0pt}{7.5mm}\textbf{Họ và tên:} \textbf{\large Nguyễn Bảo Ngọc} \hfill \textbf{Mã HS:} \texttt{BAONGOC2603} \quad \\
\hline
\end{tabularx}

\vspace{1mm}

\noindent
\begin{tabularx}{\linewidth}{|p{3.2cm}|X|}
\hline
\centering\textbf{ĐIỂM SỐ} & \centering\textbf{LỜI NHẬN XÉT CỦA THẦY CÔ} \tabularnewline
\hline
\rule{0pt}{1.3cm} & \tabularnewline
\hline
\end{tabularx}

\vspace{1.2mm}

\noindent\textbf{BẢNG ĐIỀN ĐÁP ÁN TRẮC NGHIỆM (Mỗi câu đúng 0,5 điểm):}
\vspace{0.8mm}

\noindent
\renewcommand{\arraystretch}{1.35}
\begin{tabularx}{\linewidth}{|c|*{5}{>{\centering\arraybackslash}X|}}
\hline
\textbf{Câu} & \textbf{1} & \textbf{2} & \textbf{3} & \textbf{4} & \textbf{5} \\
\hline
\textbf{Đáp án} & & & & & \\
\hline
\end{tabularx}
\renewcommand{\arraystretch}{1.0}

\vspace{1.5mm}

\noindent\textbf{\large PHẦN I. TRẮC NGHIỆM KHÁCH QUAN (2,5 điểm)}
\vspace{1mm}

\cauTN{1}{Bậc của đa thức $P = 6x^4 y^2 - 3x^3 y^3 + 2x^5 y - 6x^4 y^2 + 8$ là:}{$5$}{$6$}{$7$}{$4$}

\cauTN{2}{Đơn thức nào sau đây đồng dạng với đơn thức $-6x^3 y^3$?}{$5x^3 y^3$}{$-6x^2 y^3$}{$3x^3 y^2$}{$4x^4 y^2$}

\cauTN{3}{Kết quả của phép nhân đa thức $(3x - 2)(x + 2)$ là:}{$3x^2 + 4x - 4$}{$3x^2 - 4x - 4$}{$3x^2 + 6x - 4$}{$3x^2 - 4$}

\cauTN{4}{Kết quả của phép chia $(24x^4 y^3 - 16x^3 y^2) : 8x^2 y^2$ là:}{$3x^2 y - 2x$}{$3x^2 y - 2$}{$3x^2 y^2 - 2x$}{$3x^2 - 2y$}

\cauTN{5}{Giá trị của biểu thức $M = x^2 - 8xy + 16y^2$ tại $x = 18$ và $y = 2$ là:}{$64$}{$144$}{$100$}{$400$}

\vspace{1mm}

\noindent\textbf{\large PHẦN II. TỰ LUẬN (7,5 điểm)}

\noindent\textbf{Câu 1 (2,5 điểm). Thực hiện phép tính với đa thức:}
\vspace{0.8mm}

\noindent
\begin{tabularx}{\linewidth}{|X|X|}
\hline
\multicolumn{2}{|p{0.97\linewidth}|}{Cho hai đa thức: $A = 5x^2 y - 3xy^2 + 6xy - 3$ \ và \ $B = 2x^2 y + 3xy^2 - 4xy + 5$.} \\
\hline
\oBaiTap[7]{a)}{Tính $C = A + B$:} & 
\oBaiTap[7]{b)}{Tính $D = A - B$:} \\
\hline
\oBaiTap[5]{c)}{Tìm bậc của đa thức $C$:} & 
\oBaiTap[5]{d)}{Tính giá trị của đa thức $C$ tại $x = 1; y = -1$:} \\
\hline
\end{tabularx}


\newpage

% #############################################################################
% ========================= TOÁN 8 - TUẦN 03 - BUỔI 02: MÃ ĐỀ 103 ============
% ========================= MẶT SAU (TRANG 6) =================================
% #############################################################################

\noindent
\begin{tabularx}{\linewidth}{|X|c|}
\hline
\textbf{PHẦN II. TỰ LUẬN (Tiếp theo) - TOÁN 8 TUẦN 03 BUỔI 02} & \textbf{Mã đề: 103} \\
\hline
\end{tabularx}

\vspace{1mm}

\noindent
\begin{tabularx}{\linewidth}{|X|}
\hline
\rule[-2mm]{0pt}{7mm}\textbf{Họ và tên:} \textbf{\large Nguyễn Bảo Ngọc} \hfill \textbf{Mã HS:} \texttt{BAONGOC2603} \quad \\
\hline
\end{tabularx}

\vspace{1.5mm}

\noindent\textbf{\large Câu 2 (2,5 điểm). Rút gọn và tìm $x$:}
\vspace{0.8mm}

\noindent
\begin{tabularx}{\linewidth}{|X|X|}
\hline
\oBaiTap[12]{a)}{Rút gọn và chứng minh giá trị biểu thức không phụ thuộc vào giá trị của biến $x$:\par
$P = (2x - 3)(x + 2) - 2x(x + 1) + x + 8$.} & 
\oBaiTap[12]{b)}{Tìm số thực $x$, biết:\par
$(3x - 2)(2x + 1) - 6x(x - 1) = 23$.} \\
\hline
\end{tabularx}

\vspace{1.8mm}

\noindent\textbf{\large Câu 3 (2,5 điểm). Bài toán thực tế và nâng cao:}
\vspace{0.8mm}

\noindent
\begin{tabularx}{\linewidth}{|X|}
\hline
\vspace{0.5mm}
\textbf{a) (1,5 điểm)} Một mảnh vườn hình chữ nhật có chiều dài là $(3x + 5)\text{ m}$ và chiều rộng là $(2x + 2)\text{ m}$. Người ta làm một lối đi xung quanh vườn rộng $1\text{ m}$, phần đất còn lại bên trong dùng để trồng hoa.\par
1) Viết đa thức theo biến $x$ biểu thị diện tích mảnh vườn ban đầu.\par
2) Viết đa thức theo biến $x$ biểu thị diện tích phần đất bên trong dùng để trồng hoa.\par
\vspace{0.4mm}
\noindent\textit{Bài giải:}\dongke{9}
\vspace{1mm}
\textbf{b) (1,0 điểm)} Cho hai số $x, y$ thỏa mãn $x + y = 6$ và $xy = 5$. Không tính cụ thể giá trị của $x$ và $y$, hãy tính giá trị của biểu thức $T = x^3 + y^3$.\par
\vspace{0.4mm}
\noindent\textit{Bài giải:}\dongke{7}
\\ \hline
\end{tabularx}


\newpage

% #############################################################################
% ========================= TOÁN 8 - TUẦN 03 - BUỔI 02: MÃ ĐỀ 104 ============
% ========================= MẶT TRƯỚC (TRANG 7) ===============================
% #############################################################################

\noindent
\begin{tabularx}{\linewidth}{|X|c|}
\hline
\textbf{PHIẾU KIỂM TRA ĐỊNH KỲ TOÁN 8 - TUẦN 03 - BUỔI 02} & \textbf{Mã đề: 104} \\
\textit{Chủ đề: Ôn tập Chương 1 - Đa thức} & \textit{Thời gian: 45 phút} \\
\hline
\end{tabularx}

\vspace{1mm}

\noindent
\begin{tabularx}{\linewidth}{|X|}
\hline
\rule[-2mm]{0pt}{7.5mm}\textbf{Họ và tên:} \textbf{\large Nguyễn Bảo Ngọc} \hfill \textbf{Mã HS:} \texttt{BAONGOC0408} \quad \\
\hline
\end{tabularx}

\vspace{1mm}

\noindent
\begin{tabularx}{\linewidth}{|p{3.2cm}|X|}
\hline
\centering\textbf{ĐIỂM SỐ} & \centering\textbf{LỜI NHẬN XÉT CỦA THẦY CÔ} \tabularnewline
\hline
\rule{0pt}{1.3cm} & \tabularnewline
\hline
\end{tabularx}

\vspace{1.2mm}

\noindent\textbf{BẢNG ĐIỀN ĐÁP ÁN TRẮC NGHIỆM (Mỗi câu đúng 0,5 điểm):}
\vspace{0.8mm}

\noindent
\renewcommand{\arraystretch}{1.35}
\begin{tabularx}{\linewidth}{|c|*{5}{>{\centering\arraybackslash}X|}}
\hline
\textbf{Câu} & \textbf{1} & \textbf{2} & \textbf{3} & \textbf{4} & \textbf{5} \\
\hline
\textbf{Đáp án} & & & & & \\
\hline
\end{tabularx}
\renewcommand{\arraystretch}{1.0}

\vspace{1.5mm}

\noindent\textbf{\large PHẦN I. TRẮC NGHIỆM KHÁCH QUAN (2,5 điểm)}
\vspace{1mm}

\cauTN{1}{Bậc của đa thức $Q = 5x^4 y^3 - 2x^3 y^4 + 4x^6 y - 5x^4 y^3 + 7$ là:}{$6$}{$5$}{$7$}{$4$}

\cauTN{2}{Đơn thức nào sau đây đồng dạng với đơn thức $-7x^4 y^2$?}{$3x^4 y^2$}{$-7x^2 y^4$}{$4x^4 y^3$}{$2x^3 y^2$}

\cauTN{3}{Kết quả của phép nhân đa thức $(3x + 2)(x - 2)$ là:}{$3x^2 + 4x - 4$}{$3x^2 - 4x - 4$}{$3x^2 - 6$}{$3x^2 - 8x - 4$}

\cauTN{4}{Kết quả của phép chia $(24x^5 y^3 - 18x^3 y^2) : 6x^3 y^2$ là:}{$4x^2 y - 3$}{$4x^2 y - 3xy$}{$4x^2 y^2 - 3$}{$4x^2 - 3y$}

\cauTN{5}{Giá trị của biểu thức $N = x^2 + 8xy + 16y^2$ tại $x = 12$ và $y = 2$ là:}{$256$}{$400$}{$324$}{$100$}

\vspace{1mm}

\noindent\textbf{\large PHẦN II. TỰ LUẬN (7,5 điểm)}

\noindent\textbf{Câu 1 (2,5 điểm). Thực hiện phép tính với đa thức:}
\vspace{0.8mm}

\noindent
\begin{tabularx}{\linewidth}{|X|X|}
\hline
\multicolumn{2}{|p{0.97\linewidth}|}{Cho hai đa thức: $A = 5x^2 y - 4xy^2 + 3xy - 1$ \ và \ $B = 3x^2 y + 4xy^2 - 2xy + 6$.} \\
\hline
\oBaiTap[7]{a)}{Tính $C = A + B$:} & 
\oBaiTap[7]{b)}{Tính $D = A - B$:} \\
\hline
\oBaiTap[5]{c)}{Tìm bậc của đa thức $C$:} & 
\oBaiTap[5]{d)}{Tính giá trị của đa thức $C$ tại $x = 1; y = -1$:} \\
\hline
\end{tabularx}


\newpage

% #############################################################################
% ========================= TOÁN 8 - TUẦN 03 - BUỔI 02: MÃ ĐỀ 104 ============
% ========================= MẶT SAU (TRANG 8) =================================
% #############################################################################

\noindent
\begin{tabularx}{\linewidth}{|X|c|}
\hline
\textbf{PHẦN II. TỰ LUẬN (Tiếp theo) - TOÁN 8 TUẦN 03 BUỔI 02} & \textbf{Mã đề: 104} \\
\hline
\end{tabularx}

\vspace{1mm}

\noindent
\begin{tabularx}{\linewidth}{|X|}
\hline
\rule[-2mm]{0pt}{7mm}\textbf{Họ và tên:} \textbf{\large Nguyễn Bảo Ngọc} \hfill \textbf{Mã HS:} \texttt{BAONGOC0408} \quad \\
\hline
\end{tabularx}

\vspace{1.5mm}

\noindent\textbf{\large Câu 2 (2,5 điểm). Rút gọn và tìm $x$:}
\vspace{0.8mm}

\noindent
\begin{tabularx}{\linewidth}{|X|X|}
\hline
\oBaiTap[12]{a)}{Rút gọn và chứng minh giá trị biểu thức không phụ thuộc vào giá trị của biến $x$:\par
$N = (4x - 1)(x + 2) - 4x(x + 1) - 3x + 9$.} & 
\oBaiTap[12]{b)}{Tìm số thực $x$, biết:\par
$(2x - 1)(4x + 3) - 8x(x - 1) = 29$.} \\
\hline
\end{tabularx}

\vspace{1.8mm}

\noindent\textbf{\large Câu 3 (2,5 điểm). Bài toán thực tế và nâng cao:}
\vspace{0.8mm}

\noindent
\begin{tabularx}{\linewidth}{|X|}
\hline
\vspace{0.5mm}
\textbf{a) (1,5 điểm)} Một khu đất hình chữ nhật có chiều dài là $(4x + 5)\text{ m}$ và chiều rộng là $(2x + 4)\text{ m}$. Người ta làm một lối đi xung quanh khu đất rộng $1\text{ m}$, phần đất còn lại bên trong dùng để trồng rau sạch.\par
1) Viết đa thức theo biến $x$ biểu thị diện tích khu đất ban đầu.\par
2) Viết đa thức theo biến $x$ biểu thị diện tích phần đất bên trong dùng để trồng rau.\par
\vspace{0.4mm}
\noindent\textit{Bài giải:}\dongke{9}
\vspace{1mm}
\textbf{b) (1,0 điểm)} Cho hai số $x, y$ thỏa mãn $x + y = 6$ và $xy = 8$. Không tính cụ thể giá trị của $x$ và $y$, hãy tính giá trị của biểu thức $M = x^3 + y^3$.\par
\vspace{0.4mm}
\noindent\textit{Bài giải:}\dongke{7}
\\ \hline
\end{tabularx}



\newpage


% #############################################################################
% ========================= TOÁN 8 - TUẦN 03 - BUỔI 02: MÃ ĐỀ 101 ============
% ========================= MẶT TRƯỚC (TRANG 1) ===============================
% #############################################################################

\noindent
\begin{tabularx}{\linewidth}{|X|c|}
\hline
\textbf{PHIẾU KIỂM TRA ĐỊNH KỲ TOÁN 8 - TUẦN 03 - BUỔI 02} & \textbf{Mã đề: 101} \\
\textit{Chủ đề: Ôn tập Chương 1 - Đa thức} & \textit{Thời gian: 45 phút} \\
\hline
\end{tabularx}

\vspace{1mm}

\noindent
\begin{tabularx}{\linewidth}{|X|}
\hline
\rule[-2mm]{0pt}{7.5mm}\textbf{Họ và tên:} \textbf{\large Võ Thanh Như} \hfill \textbf{Mã HS:} \texttt{THANHNHU0404} \quad \\
\hline
\end{tabularx}

\vspace{1mm}

\noindent
\begin{tabularx}{\linewidth}{|p{3.2cm}|X|}
\hline
\centering\textbf{ĐIỂM SỐ} & \centering\textbf{LỜI NHẬN XÉT CỦA THẦY CÔ} \tabularnewline
\hline
\rule{0pt}{1.3cm} & \tabularnewline
\hline
\end{tabularx}

\vspace{1.2mm}

\noindent\textbf{BẢNG ĐIỀN ĐÁP ÁN TRẮC NGHIỆM (Mỗi câu đúng 0,5 điểm):}
\vspace{0.8mm}

\noindent
\renewcommand{\arraystretch}{1.35}
\begin{tabularx}{\linewidth}{|c|*{5}{>{\centering\arraybackslash}X|}}
\hline
\textbf{Câu} & \textbf{1} & \textbf{2} & \textbf{3} & \textbf{4} & \textbf{5} \\
\hline
\textbf{Đáp án} & & & & & \\
\hline
\end{tabularx}
\renewcommand{\arraystretch}{1.0}

\vspace{1.5mm}

\noindent\textbf{\large PHẦN I. TRẮC NGHIỆM KHÁCH QUAN (2,5 điểm)}
\vspace{1mm}

\cauTN{1}{Bậc của đa thức $M = 5x^3 y^2 - 2x^2 y^3 + 3x^4 y - 5x^3 y^2 + 1$ là:}{$4$}{$5$}{$6$}{$3$}

\cauTN{2}{Đơn thức nào sau đây đồng dạng với đơn thức $-4x^3 y^2$?}{$2x^2 y^3$}{$-4x^3 y$}{$7x^3 y^2$}{$5x^2 y^2$}

\cauTN{3}{Kết quả của phép nhân đa thức $(2x - 3)(x + 2)$ là:}{$2x^2 + x - 6$}{$2x^2 - x - 6$}{$2x^2 + 7x - 6$}{$2x^2 - 6$}

\cauTN{4}{Kết quả của phép chia $(18x^4 y^3 - 12x^3 y^2) : 6x^2 y^2$ là:}{$3x^2 y - 2xy$}{$3x^2 y - 2$}{$3x^2 y^2 - 2x$}{$3x^2 y - 2x$}

\cauTN{5}{Giá trị của biểu thức $P = x^2 - 6xy + 9y^2$ tại $x = 16$ và $y = 2$ là:}{$100$}{$144$}{$64$}{$400$}

\vspace{1mm}

\noindent\textbf{\large PHẦN II. TỰ LUẬN (7,5 điểm)}

\noindent\textbf{Câu 1 (2,5 điểm). Thực hiện phép tính với đa thức:}
\vspace{0.8mm}

\noindent
\begin{tabularx}{\linewidth}{|X|X|}
\hline
\multicolumn{2}{|p{0.97\linewidth}|}{Cho hai đa thức: $A = 3x^2 y - 2xy^2 + 5xy - 1$ \ và \ $B = 2x^2 y + 2xy^2 - 3xy + 4$.} \\
\hline
\oBaiTap[7]{a)}{Tính $C = A + B$:} & 
\oBaiTap[7]{b)}{Tính $D = A - B$:} \\
\hline
\oBaiTap[5]{c)}{Tìm bậc của đa thức $C$:} & 
\oBaiTap[5]{d)}{Tính giá trị của đa thức $C$ tại $x = 1; y = -1$:} \\
\hline
\end{tabularx}


\newpage

% #############################################################################
% ========================= TOÁN 8 - TUẦN 03 - BUỔI 02: MÃ ĐỀ 101 ============
% ========================= MẶT SAU (TRANG 2) =================================
% #############################################################################

\noindent
\begin{tabularx}{\linewidth}{|X|c|}
\hline
\textbf{PHẦN II. TỰ LUẬN (Tiếp theo) - TOÁN 8 TUẦN 03 BUỔI 02} & \textbf{Mã đề: 101} \\
\hline
\end{tabularx}

\vspace{1mm}

\noindent
\begin{tabularx}{\linewidth}{|X|}
\hline
\rule[-2mm]{0pt}{7mm}\textbf{Họ và tên:} \textbf{\large Võ Thanh Như} \hfill \textbf{Mã HS:} \texttt{THANHNHU0404} \quad \\
\hline
\end{tabularx}

\vspace{1.5mm}

\noindent\textbf{\large Câu 2 (2,5 điểm). Rút gọn và tìm $x$:}
\vspace{0.8mm}

\noindent
\begin{tabularx}{\linewidth}{|X|X|}
\hline
\oBaiTap[12]{a)}{Rút gọn và chứng minh giá trị biểu thức không phụ thuộc vào giá trị của biến $x$:\par
$M = (2x - 1)(x + 3) - 2x(x + 2) - x + 7$.} & 
\oBaiTap[12]{b)}{Tìm số thực $x$, biết:\par
$(3x - 1)(2x + 1) - 6x(x - 2) = 25$.} \\
\hline
\end{tabularx}

\vspace{1.8mm}

\noindent\textbf{\large Câu 3 (2,5 điểm). Bài toán thực tế và nâng cao:}
\vspace{0.8mm}

\noindent
\begin{tabularx}{\linewidth}{|X|}
\hline
\vspace{0.5mm}
\textbf{a) (1,5 điểm)} Một mảnh vườn hình chữ nhật có chiều dài là $(3x + 4)\text{ m}$ và chiều rộng là $(2x + 1)\text{ m}$. Người ta làm một lối đi xung quanh vườn rộng $1\text{ m}$, phần đất còn lại bên trong dùng để trồng hoa.\par
1) Viết đa thức theo biến $x$ biểu thị diện tích mảnh vườn ban đầu.\par
2) Viết đa thức theo biến $x$ biểu thị diện tích phần đất bên trong dùng để trồng hoa.\par
\vspace{0.4mm}
\noindent\textit{Bài giải:}\dongke{9}
\vspace{1mm}
\textbf{b) (1,0 điểm)} Cho hai số $x, y$ thỏa mãn $x + y = 4$ và $xy = 3$. Không tính cụ thể giá trị của $x$ và $y$, hãy tính giá trị của biểu thức $T = x^3 + y^3$.\par
\vspace{0.4mm}
\noindent\textit{Bài giải:}\dongke{7}
\\ \hline
\end{tabularx}


\newpage

% #############################################################################
% ========================= TOÁN 8 - TUẦN 03 - BUỔI 02: MÃ ĐỀ 102 ============
% ========================= MẶT TRƯỚC (TRANG 3) ===============================
% #############################################################################

\noindent
\begin{tabularx}{\linewidth}{|X|c|}
\hline
\textbf{PHIẾU KIỂM TRA ĐỊNH KỲ TOÁN 8 - TUẦN 03 - BUỔI 02} & \textbf{Mã đề: 102} \\
\textit{Chủ đề: Ôn tập Chương 1 - Đa thức} & \textit{Thời gian: 45 phút} \\
\hline
\end{tabularx}

\vspace{1mm}

\noindent
\begin{tabularx}{\linewidth}{|X|}
\hline
\rule[-2mm]{0pt}{7.5mm}\textbf{Họ và tên:} \textbf{\large Nguyễn Minh Quân} \hfill \textbf{Mã HS:} \texttt{MINHQUAN0201} \quad \\
\hline
\end{tabularx}

\vspace{1mm}

\noindent
\begin{tabularx}{\linewidth}{|p{3.2cm}|X|}
\hline
\centering\textbf{ĐIỂM SỐ} & \centering\textbf{LỜI NHẬN XÉT CỦA THẦY CÔ} \tabularnewline
\hline
\rule{0pt}{1.3cm} & \tabularnewline
\hline
\end{tabularx}

\vspace{1.2mm}

\noindent\textbf{BẢNG ĐIỀN ĐÁP ÁN TRẮC NGHIỆM (Mỗi câu đúng 0,5 điểm):}
\vspace{0.8mm}

\noindent
\renewcommand{\arraystretch}{1.35}
\begin{tabularx}{\linewidth}{|c|*{5}{>{\centering\arraybackslash}X|}}
\hline
\textbf{Câu} & \textbf{1} & \textbf{2} & \textbf{3} & \textbf{4} & \textbf{5} \\
\hline
\textbf{Đáp án} & & & & & \\
\hline
\end{tabularx}
\renewcommand{\arraystretch}{1.0}

\vspace{1.5mm}

\noindent\textbf{\large PHẦN I. TRẮC NGHIỆM KHÁCH QUAN (2,5 điểm)}
\vspace{1mm}

\cauTN{1}{Bậc của đa thức $N = 4x^3 y^3 - 3x^2 y^4 + 2x^5 y - 4x^3 y^3 + 9$ là:}{$5$}{$7$}{$6$}{$4$}

\cauTN{2}{Đơn thức nào sau đây đồng dạng với đơn thức $-5x^2 y^4$?}{$3x^2 y^4$}{$-5x^4 y^2$}{$2x^2 y^2$}{$4xy^4$}

\cauTN{3}{Kết quả của phép nhân đa thức $(2x + 3)(x - 2)$ là:}{$2x^2 + x - 6$}{$2x^2 - 7x - 6$}{$2x^2 - 6$}{$2x^2 - x - 6$}

\cauTN{4}{Kết quả của phép chia $(20x^5 y^3 - 15x^3 y^2) : 5x^3 y^2$ là:}{$4x^2 y - 3$}{$4x^2 y - 3xy$}{$4x^2 y^2 - 3$}{$4x^2 - 3y$}

\cauTN{5}{Giá trị của biểu thức $Q = x^2 + 6xy + 9y^2$ tại $x = 14$ và $y = 2$ là:}{$256$}{$400$}{$324$}{$100$}

\vspace{1mm}

\noindent\textbf{\large PHẦN II. TỰ LUẬN (7,5 điểm)}

\noindent\textbf{Câu 1 (2,5 điểm). Thực hiện phép tính với đa thức:}
\vspace{0.8mm}

\noindent
\begin{tabularx}{\linewidth}{|X|X|}
\hline
\multicolumn{2}{|p{0.97\linewidth}|}{Cho hai đa thức: $A = 4x^2 y - 3xy^2 + 4xy - 2$ \ và \ $B = 3x^2 y + 3xy^2 - 2xy + 5$.} \\
\hline
\oBaiTap[7]{a)}{Tính $C = A + B$:} & 
\oBaiTap[7]{b)}{Tính $D = A - B$:} \\
\hline
\oBaiTap[5]{c)}{Tìm bậc của đa thức $C$:} & 
\oBaiTap[5]{d)}{Tính giá trị của đa thức $C$ tại $x = 1; y = -1$:} \\
\hline
\end{tabularx}


\newpage

% #############################################################################
% ========================= TOÁN 8 - TUẦN 03 - BUỔI 02: MÃ ĐỀ 102 ============
% ========================= MẶT SAU (TRANG 4) =================================
% #############################################################################

\noindent
\begin{tabularx}{\linewidth}{|X|c|}
\hline
\textbf{PHẦN II. TỰ LUẬN (Tiếp theo) - TOÁN 8 TUẦN 03 BUỔI 02} & \textbf{Mã đề: 102} \\
\hline
\end{tabularx}

\vspace{1mm}

\noindent
\begin{tabularx}{\linewidth}{|X|}
\hline
\rule[-2mm]{0pt}{7mm}\textbf{Họ và tên:} \textbf{\large Nguyễn Minh Quân} \hfill \textbf{Mã HS:} \texttt{MINHQUAN0201} \quad \\
\hline
\end{tabularx}

\vspace{1.5mm}

\noindent\textbf{\large Câu 2 (2,5 điểm). Rút gọn và tìm $x$:}
\vspace{0.8mm}

\noindent
\begin{tabularx}{\linewidth}{|X|X|}
\hline
\oBaiTap[12]{a)}{Rút gọn và chứng minh giá trị biểu thức không phụ thuộc vào giá trị của biến $x$:\par
$N = (3x - 1)(x + 2) - 3x(x + 1) - 2x + 5$.} & 
\oBaiTap[12]{b)}{Tìm số thực $x$, biết:\par
$(2x - 1)(3x + 2) - 6x(x - 1) = 19$.} \\
\hline
\end{tabularx}

\vspace{1.8mm}

\noindent\textbf{\large Câu 3 (2,5 điểm). Bài toán thực tế và nâng cao:}
\vspace{0.8mm}

\noindent
\begin{tabularx}{\linewidth}{|X|}
\hline
\vspace{0.5mm}
\textbf{a) (1,5 điểm)} Một khu đất hình chữ nhật có chiều dài là $(4x + 3)\text{ m}$ và chiều rộng là $(2x + 3)\text{ m}$. Người ta làm một lối đi xung quanh khu đất rộng $1\text{ m}$, phần đất còn lại bên trong dùng để trồng rau sạch.\par
1) Viết đa thức theo biến $x$ biểu thị diện tích khu đất ban đầu.\par
2) Viết đa thức theo biến $x$ biểu thị diện tích phần đất bên trong dùng để trồng rau.\par
\vspace{0.4mm}
\noindent\textit{Bài giải:}\dongke{9}
\vspace{1mm}
\textbf{b) (1,0 điểm)} Cho hai số $x, y$ thỏa mãn $x + y = 5$ và $xy = 4$. Không tính cụ thể giá trị của $x$ và $y$, hãy tính giá trị của biểu thức $M = x^3 + y^3$.\par
\vspace{0.4mm}
\noindent\textit{Bài giải:}\dongke{7}
\\ \hline
\end{tabularx}


\newpage

% #############################################################################
% ========================= TOÁN 8 - TUẦN 03 - BUỔI 02: MÃ ĐỀ 103 ============
% ========================= MẶT TRƯỚC (TRANG 5) ===============================
% #############################################################################

\noindent
\begin{tabularx}{\linewidth}{|X|c|}
\hline
\textbf{PHIẾU KIỂM TRA ĐỊNH KỲ TOÁN 8 - TUẦN 03 - BUỔI 02} & \textbf{Mã đề: 103} \\
\textit{Chủ đề: Ôn tập Chương 1 - Đa thức} & \textit{Thời gian: 45 phút} \\
\hline
\end{tabularx}

\vspace{1mm}

\noindent
\begin{tabularx}{\linewidth}{|X|}
\hline
\rule[-2mm]{0pt}{7.5mm}\textbf{Họ và tên:} \textbf{\large Lưu Thị Thùy Trang} \hfill \textbf{Mã HS:} \texttt{THUYTRANG0305} \quad \\
\hline
\end{tabularx}

\vspace{1mm}

\noindent
\begin{tabularx}{\linewidth}{|p{3.2cm}|X|}
\hline
\centering\textbf{ĐIỂM SỐ} & \centering\textbf{LỜI NHẬN XÉT CỦA THẦY CÔ} \tabularnewline
\hline
\rule{0pt}{1.3cm} & \tabularnewline
\hline
\end{tabularx}

\vspace{1.2mm}

\noindent\textbf{BẢNG ĐIỀN ĐÁP ÁN TRẮC NGHIỆM (Mỗi câu đúng 0,5 điểm):}
\vspace{0.8mm}

\noindent
\renewcommand{\arraystretch}{1.35}
\begin{tabularx}{\linewidth}{|c|*{5}{>{\centering\arraybackslash}X|}}
\hline
\textbf{Câu} & \textbf{1} & \textbf{2} & \textbf{3} & \textbf{4} & \textbf{5} \\
\hline
\textbf{Đáp án} & & & & & \\
\hline
\end{tabularx}
\renewcommand{\arraystretch}{1.0}

\vspace{1.5mm}

\noindent\textbf{\large PHẦN I. TRẮC NGHIỆM KHÁCH QUAN (2,5 điểm)}
\vspace{1mm}

\cauTN{1}{Bậc của đa thức $P = 6x^4 y^2 - 3x^3 y^3 + 2x^5 y - 6x^4 y^2 + 8$ là:}{$5$}{$6$}{$7$}{$4$}

\cauTN{2}{Đơn thức nào sau đây đồng dạng với đơn thức $-6x^3 y^3$?}{$5x^3 y^3$}{$-6x^2 y^3$}{$3x^3 y^2$}{$4x^4 y^2$}

\cauTN{3}{Kết quả của phép nhân đa thức $(3x - 2)(x + 2)$ là:}{$3x^2 + 4x - 4$}{$3x^2 - 4x - 4$}{$3x^2 + 6x - 4$}{$3x^2 - 4$}

\cauTN{4}{Kết quả của phép chia $(24x^4 y^3 - 16x^3 y^2) : 8x^2 y^2$ là:}{$3x^2 y - 2x$}{$3x^2 y - 2$}{$3x^2 y^2 - 2x$}{$3x^2 - 2y$}

\cauTN{5}{Giá trị của biểu thức $M = x^2 - 8xy + 16y^2$ tại $x = 18$ và $y = 2$ là:}{$64$}{$144$}{$100$}{$400$}

\vspace{1mm}

\noindent\textbf{\large PHẦN II. TỰ LUẬN (7,5 điểm)}

\noindent\textbf{Câu 1 (2,5 điểm). Thực hiện phép tính với đa thức:}
\vspace{0.8mm}

\noindent
\begin{tabularx}{\linewidth}{|X|X|}
\hline
\multicolumn{2}{|p{0.97\linewidth}|}{Cho hai đa thức: $A = 5x^2 y - 3xy^2 + 6xy - 3$ \ và \ $B = 2x^2 y + 3xy^2 - 4xy + 5$.} \\
\hline
\oBaiTap[7]{a)}{Tính $C = A + B$:} & 
\oBaiTap[7]{b)}{Tính $D = A - B$:} \\
\hline
\oBaiTap[5]{c)}{Tìm bậc của đa thức $C$:} & 
\oBaiTap[5]{d)}{Tính giá trị của đa thức $C$ tại $x = 1; y = -1$:} \\
\hline
\end{tabularx}


\newpage

% #############################################################################
% ========================= TOÁN 8 - TUẦN 03 - BUỔI 02: MÃ ĐỀ 103 ============
% ========================= MẶT SAU (TRANG 6) =================================
% #############################################################################

\noindent
\begin{tabularx}{\linewidth}{|X|c|}
\hline
\textbf{PHẦN II. TỰ LUẬN (Tiếp theo) - TOÁN 8 TUẦN 03 BUỔI 02} & \textbf{Mã đề: 103} \\
\hline
\end{tabularx}

\vspace{1mm}

\noindent
\begin{tabularx}{\linewidth}{|X|}
\hline
\rule[-2mm]{0pt}{7mm}\textbf{Họ và tên:} \textbf{\large Lưu Thị Thùy Trang} \hfill \textbf{Mã HS:} \texttt{THUYTRANG0305} \quad \\
\hline
\end{tabularx}

\vspace{1.5mm}

\noindent\textbf{\large Câu 2 (2,5 điểm). Rút gọn và tìm $x$:}
\vspace{0.8mm}

\noindent
\begin{tabularx}{\linewidth}{|X|X|}
\hline
\oBaiTap[12]{a)}{Rút gọn và chứng minh giá trị biểu thức không phụ thuộc vào giá trị của biến $x$:\par
$P = (2x - 3)(x + 2) - 2x(x + 1) + x + 8$.} & 
\oBaiTap[12]{b)}{Tìm số thực $x$, biết:\par
$(3x - 2)(2x + 1) - 6x(x - 1) = 23$.} \\
\hline
\end{tabularx}

\vspace{1.8mm}

\noindent\textbf{\large Câu 3 (2,5 điểm). Bài toán thực tế và nâng cao:}
\vspace{0.8mm}

\noindent
\begin{tabularx}{\linewidth}{|X|}
\hline
\vspace{0.5mm}
\textbf{a) (1,5 điểm)} Một mảnh vườn hình chữ nhật có chiều dài là $(3x + 5)\text{ m}$ và chiều rộng là $(2x + 2)\text{ m}$. Người ta làm một lối đi xung quanh vườn rộng $1\text{ m}$, phần đất còn lại bên trong dùng để trồng hoa.\par
1) Viết đa thức theo biến $x$ biểu thị diện tích mảnh vườn ban đầu.\par
2) Viết đa thức theo biến $x$ biểu thị diện tích phần đất bên trong dùng để trồng hoa.\par
\vspace{0.4mm}
\noindent\textit{Bài giải:}\dongke{9}
\vspace{1mm}
\textbf{b) (1,0 điểm)} Cho hai số $x, y$ thỏa mãn $x + y = 6$ và $xy = 5$. Không tính cụ thể giá trị của $x$ và $y$, hãy tính giá trị của biểu thức $T = x^3 + y^3$.\par
\vspace{0.4mm}
\noindent\textit{Bài giải:}\dongke{7}
\\ \hline
\end{tabularx}


\newpage

% #############################################################################
% ========================= TOÁN 8 - TUẦN 03 - BUỔI 02: MÃ ĐỀ 104 ============
% ========================= MẶT TRƯỚC (TRANG 7) ===============================
% #############################################################################

\noindent
\begin{tabularx}{\linewidth}{|X|c|}
\hline
\textbf{PHIẾU KIỂM TRA ĐỊNH KỲ TOÁN 8 - TUẦN 03 - BUỔI 02} & \textbf{Mã đề: 104} \\
\textit{Chủ đề: Ôn tập Chương 1 - Đa thức} & \textit{Thời gian: 45 phút} \\
\hline
\end{tabularx}

\vspace{1mm}

\noindent
\begin{tabularx}{\linewidth}{|X|}
\hline
\rule[-2mm]{0pt}{7.5mm}\textbf{Họ và tên:} \textbf{\large Trần Thế Vinh} \hfill \textbf{Mã HS:} \texttt{THEVINH0507} \quad \\
\hline
\end{tabularx}

\vspace{1mm}

\noindent
\begin{tabularx}{\linewidth}{|p{3.2cm}|X|}
\hline
\centering\textbf{ĐIỂM SỐ} & \centering\textbf{LỜI NHẬN XÉT CỦA THẦY CÔ} \tabularnewline
\hline
\rule{0pt}{1.3cm} & \tabularnewline
\hline
\end{tabularx}

\vspace{1.2mm}

\noindent\textbf{BẢNG ĐIỀN ĐÁP ÁN TRẮC NGHIỆM (Mỗi câu đúng 0,5 điểm):}
\vspace{0.8mm}

\noindent
\renewcommand{\arraystretch}{1.35}
\begin{tabularx}{\linewidth}{|c|*{5}{>{\centering\arraybackslash}X|}}
\hline
\textbf{Câu} & \textbf{1} & \textbf{2} & \textbf{3} & \textbf{4} & \textbf{5} \\
\hline
\textbf{Đáp án} & & & & & \\
\hline
\end{tabularx}
\renewcommand{\arraystretch}{1.0}

\vspace{1.5mm}

\noindent\textbf{\large PHẦN I. TRẮC NGHIỆM KHÁCH QUAN (2,5 điểm)}
\vspace{1mm}

\cauTN{1}{Bậc của đa thức $Q = 5x^4 y^3 - 2x^3 y^4 + 4x^6 y - 5x^4 y^3 + 7$ là:}{$6$}{$5$}{$7$}{$4$}

\cauTN{2}{Đơn thức nào sau đây đồng dạng với đơn thức $-7x^4 y^2$?}{$3x^4 y^2$}{$-7x^2 y^4$}{$4x^4 y^3$}{$2x^3 y^2$}

\cauTN{3}{Kết quả của phép nhân đa thức $(3x + 2)(x - 2)$ là:}{$3x^2 + 4x - 4$}{$3x^2 - 4x - 4$}{$3x^2 - 6$}{$3x^2 - 8x - 4$}

\cauTN{4}{Kết quả của phép chia $(24x^5 y^3 - 18x^3 y^2) : 6x^3 y^2$ là:}{$4x^2 y - 3$}{$4x^2 y - 3xy$}{$4x^2 y^2 - 3$}{$4x^2 - 3y$}

\cauTN{5}{Giá trị của biểu thức $N = x^2 + 8xy + 16y^2$ tại $x = 12$ và $y = 2$ là:}{$256$}{$400$}{$324$}{$100$}

\vspace{1mm}

\noindent\textbf{\large PHẦN II. TỰ LUẬN (7,5 điểm)}

\noindent\textbf{Câu 1 (2,5 điểm). Thực hiện phép tính với đa thức:}
\vspace{0.8mm}

\noindent
\begin{tabularx}{\linewidth}{|X|X|}
\hline
\multicolumn{2}{|p{0.97\linewidth}|}{Cho hai đa thức: $A = 5x^2 y - 4xy^2 + 3xy - 1$ \ và \ $B = 3x^2 y + 4xy^2 - 2xy + 6$.} \\
\hline
\oBaiTap[7]{a)}{Tính $C = A + B$:} & 
\oBaiTap[7]{b)}{Tính $D = A - B$:} \\
\hline
\oBaiTap[5]{c)}{Tìm bậc của đa thức $C$:} & 
\oBaiTap[5]{d)}{Tính giá trị của đa thức $C$ tại $x = 1; y = -1$:} \\
\hline
\end{tabularx}


\newpage

% #############################################################################
% ========================= TOÁN 8 - TUẦN 03 - BUỔI 02: MÃ ĐỀ 104 ============
% ========================= MẶT SAU (TRANG 8) =================================
% #############################################################################

\noindent
\begin{tabularx}{\linewidth}{|X|c|}
\hline
\textbf{PHẦN II. TỰ LUẬN (Tiếp theo) - TOÁN 8 TUẦN 03 BUỔI 02} & \textbf{Mã đề: 104} \\
\hline
\end{tabularx}

\vspace{1mm}

\noindent
\begin{tabularx}{\linewidth}{|X|}
\hline
\rule[-2mm]{0pt}{7mm}\textbf{Họ và tên:} \textbf{\large Trần Thế Vinh} \hfill \textbf{Mã HS:} \texttt{THEVINH0507} \quad \\
\hline
\end{tabularx}

\vspace{1.5mm}

\noindent\textbf{\large Câu 2 (2,5 điểm). Rút gọn và tìm $x$:}
\vspace{0.8mm}

\noindent
\begin{tabularx}{\linewidth}{|X|X|}
\hline
\oBaiTap[12]{a)}{Rút gọn và chứng minh giá trị biểu thức không phụ thuộc vào giá trị của biến $x$:\par
$N = (4x - 1)(x + 2) - 4x(x + 1) - 3x + 9$.} & 
\oBaiTap[12]{b)}{Tìm số thực $x$, biết:\par
$(2x - 1)(4x + 3) - 8x(x - 1) = 29$.} \\
\hline
\end{tabularx}

\vspace{1.8mm}

\noindent\textbf{\large Câu 3 (2,5 điểm). Bài toán thực tế và nâng cao:}
\vspace{0.8mm}

\noindent
\begin{tabularx}{\linewidth}{|X|}
\hline
\vspace{0.5mm}
\textbf{a) (1,5 điểm)} Một khu đất hình chữ nhật có chiều dài là $(4x + 5)\text{ m}$ và chiều rộng là $(2x + 4)\text{ m}$. Người ta làm một lối đi xung quanh khu đất rộng $1\text{ m}$, phần đất còn lại bên trong dùng để trồng rau sạch.\par
1) Viết đa thức theo biến $x$ biểu thị diện tích khu đất ban đầu.\par
2) Viết đa thức theo biến $x$ biểu thị diện tích phần đất bên trong dùng để trồng rau.\par
\vspace{0.4mm}
\noindent\textit{Bài giải:}\dongke{9}
\vspace{1mm}
\textbf{b) (1,0 điểm)} Cho hai số $x, y$ thỏa mãn $x + y = 6$ và $xy = 8$. Không tính cụ thể giá trị của $x$ và $y$, hãy tính giá trị của biểu thức $M = x^3 + y^3$.\par
\vspace{0.4mm}
\noindent\textit{Bài giải:}\dongke{7}
\\ \hline
\end{tabularx}



\newpage


% #############################################################################
% ========================= TOÁN 8 - TUẦN 03 - BUỔI 02: MÃ ĐỀ 101 ============
% ========================= MẶT TRƯỚC (TRANG 1) ===============================
% #############################################################################

\noindent
\begin{tabularx}{\linewidth}{|X|c|}
\hline
\textbf{PHIẾU KIỂM TRA ĐỊNH KỲ TOÁN 8 - TUẦN 03 - BUỔI 02} & \textbf{Mã đề: 101} \\
\textit{Chủ đề: Ôn tập Chương 1 - Đa thức} & \textit{Thời gian: 45 phút} \\
\hline
\end{tabularx}

\vspace{1mm}

\noindent
\begin{tabularx}{\linewidth}{|X|}
\hline
\rule[-2mm]{0pt}{7.5mm}\textbf{Họ và tên:} \textbf{\large Hoàng Hà Vy} \hfill \textbf{Mã HS:} \texttt{HAVY0403} \quad \\
\hline
\end{tabularx}

\vspace{1mm}

\noindent
\begin{tabularx}{\linewidth}{|p{3.2cm}|X|}
\hline
\centering\textbf{ĐIỂM SỐ} & \centering\textbf{LỜI NHẬN XÉT CỦA THẦY CÔ} \tabularnewline
\hline
\rule{0pt}{1.3cm} & \tabularnewline
\hline
\end{tabularx}

\vspace{1.2mm}

\noindent\textbf{BẢNG ĐIỀN ĐÁP ÁN TRẮC NGHIỆM (Mỗi câu đúng 0,5 điểm):}
\vspace{0.8mm}

\noindent
\renewcommand{\arraystretch}{1.35}
\begin{tabularx}{\linewidth}{|c|*{5}{>{\centering\arraybackslash}X|}}
\hline
\textbf{Câu} & \textbf{1} & \textbf{2} & \textbf{3} & \textbf{4} & \textbf{5} \\
\hline
\textbf{Đáp án} & & & & & \\
\hline
\end{tabularx}
\renewcommand{\arraystretch}{1.0}

\vspace{1.5mm}

\noindent\textbf{\large PHẦN I. TRẮC NGHIỆM KHÁCH QUAN (2,5 điểm)}
\vspace{1mm}

\cauTN{1}{Bậc của đa thức $M = 5x^3 y^2 - 2x^2 y^3 + 3x^4 y - 5x^3 y^2 + 1$ là:}{$4$}{$5$}{$6$}{$3$}

\cauTN{2}{Đơn thức nào sau đây đồng dạng với đơn thức $-4x^3 y^2$?}{$2x^2 y^3$}{$-4x^3 y$}{$7x^3 y^2$}{$5x^2 y^2$}

\cauTN{3}{Kết quả của phép nhân đa thức $(2x - 3)(x + 2)$ là:}{$2x^2 + x - 6$}{$2x^2 - x - 6$}{$2x^2 + 7x - 6$}{$2x^2 - 6$}

\cauTN{4}{Kết quả của phép chia $(18x^4 y^3 - 12x^3 y^2) : 6x^2 y^2$ là:}{$3x^2 y - 2xy$}{$3x^2 y - 2$}{$3x^2 y^2 - 2x$}{$3x^2 y - 2x$}

\cauTN{5}{Giá trị của biểu thức $P = x^2 - 6xy + 9y^2$ tại $x = 16$ và $y = 2$ là:}{$100$}{$144$}{$64$}{$400$}

\vspace{1mm}

\noindent\textbf{\large PHẦN II. TỰ LUẬN (7,5 điểm)}

\noindent\textbf{Câu 1 (2,5 điểm). Thực hiện phép tính với đa thức:}
\vspace{0.8mm}

\noindent
\begin{tabularx}{\linewidth}{|X|X|}
\hline
\multicolumn{2}{|p{0.97\linewidth}|}{Cho hai đa thức: $A = 3x^2 y - 2xy^2 + 5xy - 1$ \ và \ $B = 2x^2 y + 2xy^2 - 3xy + 4$.} \\
\hline
\oBaiTap[7]{a)}{Tính $C = A + B$:} & 
\oBaiTap[7]{b)}{Tính $D = A - B$:} \\
\hline
\oBaiTap[5]{c)}{Tìm bậc của đa thức $C$:} & 
\oBaiTap[5]{d)}{Tính giá trị của đa thức $C$ tại $x = 1; y = -1$:} \\
\hline
\end{tabularx}


\newpage

% #############################################################################
% ========================= TOÁN 8 - TUẦN 03 - BUỔI 02: MÃ ĐỀ 101 ============
% ========================= MẶT SAU (TRANG 2) =================================
% #############################################################################

\noindent
\begin{tabularx}{\linewidth}{|X|c|}
\hline
\textbf{PHẦN II. TỰ LUẬN (Tiếp theo) - TOÁN 8 TUẦN 03 BUỔI 02} & \textbf{Mã đề: 101} \\
\hline
\end{tabularx}

\vspace{1mm}

\noindent
\begin{tabularx}{\linewidth}{|X|}
\hline
\rule[-2mm]{0pt}{7mm}\textbf{Họ và tên:} \textbf{\large Hoàng Hà Vy} \hfill \textbf{Mã HS:} \texttt{HAVY0403} \quad \\
\hline
\end{tabularx}

\vspace{1.5mm}

\noindent\textbf{\large Câu 2 (2,5 điểm). Rút gọn và tìm $x$:}
\vspace{0.8mm}

\noindent
\begin{tabularx}{\linewidth}{|X|X|}
\hline
\oBaiTap[12]{a)}{Rút gọn và chứng minh giá trị biểu thức không phụ thuộc vào giá trị của biến $x$:\par
$M = (2x - 1)(x + 3) - 2x(x + 2) - x + 7$.} & 
\oBaiTap[12]{b)}{Tìm số thực $x$, biết:\par
$(3x - 1)(2x + 1) - 6x(x - 2) = 25$.} \\
\hline
\end{tabularx}

\vspace{1.8mm}

\noindent\textbf{\large Câu 3 (2,5 điểm). Bài toán thực tế và nâng cao:}
\vspace{0.8mm}

\noindent
\begin{tabularx}{\linewidth}{|X|}
\hline
\vspace{0.5mm}
\textbf{a) (1,5 điểm)} Một mảnh vườn hình chữ nhật có chiều dài là $(3x + 4)\text{ m}$ và chiều rộng là $(2x + 1)\text{ m}$. Người ta làm một lối đi xung quanh vườn rộng $1\text{ m}$, phần đất còn lại bên trong dùng để trồng hoa.\par
1) Viết đa thức theo biến $x$ biểu thị diện tích mảnh vườn ban đầu.\par
2) Viết đa thức theo biến $x$ biểu thị diện tích phần đất bên trong dùng để trồng hoa.\par
\vspace{0.4mm}
\noindent\textit{Bài giải:}\dongke{9}
\vspace{1mm}
\textbf{b) (1,0 điểm)} Cho hai số $x, y$ thỏa mãn $x + y = 4$ và $xy = 3$. Không tính cụ thể giá trị của $x$ và $y$, hãy tính giá trị của biểu thức $T = x^3 + y^3$.\par
\vspace{0.4mm}
\noindent\textit{Bài giải:}\dongke{7}
\\ \hline
\end{tabularx}


\newpage

% #############################################################################
% ========================= TOÁN 8 - TUẦN 03 - BUỔI 02: MÃ ĐỀ 102 ============
% ========================= MẶT TRƯỚC (TRANG 3) ===============================
% #############################################################################

\noindent
\begin{tabularx}{\linewidth}{|X|c|}
\hline
\textbf{PHIẾU KIỂM TRA ĐỊNH KỲ TOÁN 8 - TUẦN 03 - BUỔI 02} & \textbf{Mã đề: 102} \\
\textit{Chủ đề: Ôn tập Chương 1 - Đa thức} & \textit{Thời gian: 45 phút} \\
\hline
\end{tabularx}

\vspace{1mm}

\noindent
\begin{tabularx}{\linewidth}{|X|}
\hline
\rule[-2mm]{0pt}{7.5mm}\textbf{Họ và tên:} \textbf{\large Trương Bảo Vy} \hfill \textbf{Mã HS:} \texttt{BAOVY3006} \quad \\
\hline
\end{tabularx}

\vspace{1mm}

\noindent
\begin{tabularx}{\linewidth}{|p{3.2cm}|X|}
\hline
\centering\textbf{ĐIỂM SỐ} & \centering\textbf{LỜI NHẬN XÉT CỦA THẦY CÔ} \tabularnewline
\hline
\rule{0pt}{1.3cm} & \tabularnewline
\hline
\end{tabularx}

\vspace{1.2mm}

\noindent\textbf{BẢNG ĐIỀN ĐÁP ÁN TRẮC NGHIỆM (Mỗi câu đúng 0,5 điểm):}
\vspace{0.8mm}

\noindent
\renewcommand{\arraystretch}{1.35}
\begin{tabularx}{\linewidth}{|c|*{5}{>{\centering\arraybackslash}X|}}
\hline
\textbf{Câu} & \textbf{1} & \textbf{2} & \textbf{3} & \textbf{4} & \textbf{5} \\
\hline
\textbf{Đáp án} & & & & & \\
\hline
\end{tabularx}
\renewcommand{\arraystretch}{1.0}

\vspace{1.5mm}

\noindent\textbf{\large PHẦN I. TRẮC NGHIỆM KHÁCH QUAN (2,5 điểm)}
\vspace{1mm}

\cauTN{1}{Bậc của đa thức $N = 4x^3 y^3 - 3x^2 y^4 + 2x^5 y - 4x^3 y^3 + 9$ là:}{$5$}{$7$}{$6$}{$4$}

\cauTN{2}{Đơn thức nào sau đây đồng dạng với đơn thức $-5x^2 y^4$?}{$3x^2 y^4$}{$-5x^4 y^2$}{$2x^2 y^2$}{$4xy^4$}

\cauTN{3}{Kết quả của phép nhân đa thức $(2x + 3)(x - 2)$ là:}{$2x^2 + x - 6$}{$2x^2 - 7x - 6$}{$2x^2 - 6$}{$2x^2 - x - 6$}

\cauTN{4}{Kết quả của phép chia $(20x^5 y^3 - 15x^3 y^2) : 5x^3 y^2$ là:}{$4x^2 y - 3$}{$4x^2 y - 3xy$}{$4x^2 y^2 - 3$}{$4x^2 - 3y$}

\cauTN{5}{Giá trị của biểu thức $Q = x^2 + 6xy + 9y^2$ tại $x = 14$ và $y = 2$ là:}{$256$}{$400$}{$324$}{$100$}

\vspace{1mm}

\noindent\textbf{\large PHẦN II. TỰ LUẬN (7,5 điểm)}

\noindent\textbf{Câu 1 (2,5 điểm). Thực hiện phép tính với đa thức:}
\vspace{0.8mm}

\noindent
\begin{tabularx}{\linewidth}{|X|X|}
\hline
\multicolumn{2}{|p{0.97\linewidth}|}{Cho hai đa thức: $A = 4x^2 y - 3xy^2 + 4xy - 2$ \ và \ $B = 3x^2 y + 3xy^2 - 2xy + 5$.} \\
\hline
\oBaiTap[7]{a)}{Tính $C = A + B$:} & 
\oBaiTap[7]{b)}{Tính $D = A - B$:} \\
\hline
\oBaiTap[5]{c)}{Tìm bậc của đa thức $C$:} & 
\oBaiTap[5]{d)}{Tính giá trị của đa thức $C$ tại $x = 1; y = -1$:} \\
\hline
\end{tabularx}


\newpage

% #############################################################################
% ========================= TOÁN 8 - TUẦN 03 - BUỔI 02: MÃ ĐỀ 102 ============
% ========================= MẶT SAU (TRANG 4) =================================
% #############################################################################

\noindent
\begin{tabularx}{\linewidth}{|X|c|}
\hline
\textbf{PHẦN II. TỰ LUẬN (Tiếp theo) - TOÁN 8 TUẦN 03 BUỔI 02} & \textbf{Mã đề: 102} \\
\hline
\end{tabularx}

\vspace{1mm}

\noindent
\begin{tabularx}{\linewidth}{|X|}
\hline
\rule[-2mm]{0pt}{7mm}\textbf{Họ và tên:} \textbf{\large Trương Bảo Vy} \hfill \textbf{Mã HS:} \texttt{BAOVY3006} \quad \\
\hline
\end{tabularx}

\vspace{1.5mm}

\noindent\textbf{\large Câu 2 (2,5 điểm). Rút gọn và tìm $x$:}
\vspace{0.8mm}

\noindent
\begin{tabularx}{\linewidth}{|X|X|}
\hline
\oBaiTap[12]{a)}{Rút gọn và chứng minh giá trị biểu thức không phụ thuộc vào giá trị của biến $x$:\par
$N = (3x - 1)(x + 2) - 3x(x + 1) - 2x + 5$.} & 
\oBaiTap[12]{b)}{Tìm số thực $x$, biết:\par
$(2x - 1)(3x + 2) - 6x(x - 1) = 19$.} \\
\hline
\end{tabularx}

\vspace{1.8mm}

\noindent\textbf{\large Câu 3 (2,5 điểm). Bài toán thực tế và nâng cao:}
\vspace{0.8mm}

\noindent
\begin{tabularx}{\linewidth}{|X|}
\hline
\vspace{0.5mm}
\textbf{a) (1,5 điểm)} Một khu đất hình chữ nhật có chiều dài là $(4x + 3)\text{ m}$ và chiều rộng là $(2x + 3)\text{ m}$. Người ta làm một lối đi xung quanh khu đất rộng $1\text{ m}$, phần đất còn lại bên trong dùng để trồng rau sạch.\par
1) Viết đa thức theo biến $x$ biểu thị diện tích khu đất ban đầu.\par
2) Viết đa thức theo biến $x$ biểu thị diện tích phần đất bên trong dùng để trồng rau.\par
\vspace{0.4mm}
\noindent\textit{Bài giải:}\dongke{9}
\vspace{1mm}
\textbf{b) (1,0 điểm)} Cho hai số $x, y$ thỏa mãn $x + y = 5$ và $xy = 4$. Không tính cụ thể giá trị của $x$ và $y$, hãy tính giá trị của biểu thức $M = x^3 + y^3$.\par
\vspace{0.4mm}
\noindent\textit{Bài giải:}\dongke{7}
\\ \hline
\end{tabularx}


\newpage

% #############################################################################
% ========================= TOÁN 8 - TUẦN 03 - BUỔI 02: MÃ ĐỀ 103 ============
% ========================= MẶT TRƯỚC (TRANG 5) ===============================
% #############################################################################

\noindent
\begin{tabularx}{\linewidth}{|X|c|}
\hline
\textbf{PHIẾU KIỂM TRA ĐỊNH KỲ TOÁN 8 - TUẦN 03 - BUỔI 02} & \textbf{Mã đề: 103} \\
\textit{Chủ đề: Ôn tập Chương 1 - Đa thức} & \textit{Thời gian: 45 phút} \\
\hline
\end{tabularx}

\vspace{1mm}

\noindent
\begin{tabularx}{\linewidth}{|X|}
\hline
\rule[-2mm]{0pt}{7.5mm}\textbf{Họ và tên:} \textbf{\large Đào Nguyễn Kim Ánh} \hfill \textbf{Mã HS:} \texttt{KIMANH0804} \quad \\
\hline
\end{tabularx}

\vspace{1mm}

\noindent
\begin{tabularx}{\linewidth}{|p{3.2cm}|X|}
\hline
\centering\textbf{ĐIỂM SỐ} & \centering\textbf{LỜI NHẬN XÉT CỦA THẦY CÔ} \tabularnewline
\hline
\rule{0pt}{1.3cm} & \tabularnewline
\hline
\end{tabularx}

\vspace{1.2mm}

\noindent\textbf{BẢNG ĐIỀN ĐÁP ÁN TRẮC NGHIỆM (Mỗi câu đúng 0,5 điểm):}
\vspace{0.8mm}

\noindent
\renewcommand{\arraystretch}{1.35}
\begin{tabularx}{\linewidth}{|c|*{5}{>{\centering\arraybackslash}X|}}
\hline
\textbf{Câu} & \textbf{1} & \textbf{2} & \textbf{3} & \textbf{4} & \textbf{5} \\
\hline
\textbf{Đáp án} & & & & & \\
\hline
\end{tabularx}
\renewcommand{\arraystretch}{1.0}

\vspace{1.5mm}

\noindent\textbf{\large PHẦN I. TRẮC NGHIỆM KHÁCH QUAN (2,5 điểm)}
\vspace{1mm}

\cauTN{1}{Bậc của đa thức $P = 6x^4 y^2 - 3x^3 y^3 + 2x^5 y - 6x^4 y^2 + 8$ là:}{$5$}{$6$}{$7$}{$4$}

\cauTN{2}{Đơn thức nào sau đây đồng dạng với đơn thức $-6x^3 y^3$?}{$5x^3 y^3$}{$-6x^2 y^3$}{$3x^3 y^2$}{$4x^4 y^2$}

\cauTN{3}{Kết quả của phép nhân đa thức $(3x - 2)(x + 2)$ là:}{$3x^2 + 4x - 4$}{$3x^2 - 4x - 4$}{$3x^2 + 6x - 4$}{$3x^2 - 4$}

\cauTN{4}{Kết quả của phép chia $(24x^4 y^3 - 16x^3 y^2) : 8x^2 y^2$ là:}{$3x^2 y - 2x$}{$3x^2 y - 2$}{$3x^2 y^2 - 2x$}{$3x^2 - 2y$}

\cauTN{5}{Giá trị của biểu thức $M = x^2 - 8xy + 16y^2$ tại $x = 18$ và $y = 2$ là:}{$64$}{$144$}{$100$}{$400$}

\vspace{1mm}

\noindent\textbf{\large PHẦN II. TỰ LUẬN (7,5 điểm)}

\noindent\textbf{Câu 1 (2,5 điểm). Thực hiện phép tính với đa thức:}
\vspace{0.8mm}

\noindent
\begin{tabularx}{\linewidth}{|X|X|}
\hline
\multicolumn{2}{|p{0.97\linewidth}|}{Cho hai đa thức: $A = 5x^2 y - 3xy^2 + 6xy - 3$ \ và \ $B = 2x^2 y + 3xy^2 - 4xy + 5$.} \\
\hline
\oBaiTap[7]{a)}{Tính $C = A + B$:} & 
\oBaiTap[7]{b)}{Tính $D = A - B$:} \\
\hline
\oBaiTap[5]{c)}{Tìm bậc của đa thức $C$:} & 
\oBaiTap[5]{d)}{Tính giá trị của đa thức $C$ tại $x = 1; y = -1$:} \\
\hline
\end{tabularx}


\newpage

% #############################################################################
% ========================= TOÁN 8 - TUẦN 03 - BUỔI 02: MÃ ĐỀ 103 ============
% ========================= MẶT SAU (TRANG 6) =================================
% #############################################################################

\noindent
\begin{tabularx}{\linewidth}{|X|c|}
\hline
\textbf{PHẦN II. TỰ LUẬN (Tiếp theo) - TOÁN 8 TUẦN 03 BUỔI 02} & \textbf{Mã đề: 103} \\
\hline
\end{tabularx}

\vspace{1mm}

\noindent
\begin{tabularx}{\linewidth}{|X|}
\hline
\rule[-2mm]{0pt}{7mm}\textbf{Họ và tên:} \textbf{\large Đào Nguyễn Kim Ánh} \hfill \textbf{Mã HS:} \texttt{KIMANH0804} \quad \\
\hline
\end{tabularx}

\vspace{1.5mm}

\noindent\textbf{\large Câu 2 (2,5 điểm). Rút gọn và tìm $x$:}
\vspace{0.8mm}

\noindent
\begin{tabularx}{\linewidth}{|X|X|}
\hline
\oBaiTap[12]{a)}{Rút gọn và chứng minh giá trị biểu thức không phụ thuộc vào giá trị của biến $x$:\par
$P = (2x - 3)(x + 2) - 2x(x + 1) + x + 8$.} & 
\oBaiTap[12]{b)}{Tìm số thực $x$, biết:\par
$(3x - 2)(2x + 1) - 6x(x - 1) = 23$.} \\
\hline
\end{tabularx}

\vspace{1.8mm}

\noindent\textbf{\large Câu 3 (2,5 điểm). Bài toán thực tế và nâng cao:}
\vspace{0.8mm}

\noindent
\begin{tabularx}{\linewidth}{|X|}
\hline
\vspace{0.5mm}
\textbf{a) (1,5 điểm)} Một mảnh vườn hình chữ nhật có chiều dài là $(3x + 5)\text{ m}$ và chiều rộng là $(2x + 2)\text{ m}$. Người ta làm một lối đi xung quanh vườn rộng $1\text{ m}$, phần đất còn lại bên trong dùng để trồng hoa.\par
1) Viết đa thức theo biến $x$ biểu thị diện tích mảnh vườn ban đầu.\par
2) Viết đa thức theo biến $x$ biểu thị diện tích phần đất bên trong dùng để trồng hoa.\par
\vspace{0.4mm}
\noindent\textit{Bài giải:}\dongke{9}
\vspace{1mm}
\textbf{b) (1,0 điểm)} Cho hai số $x, y$ thỏa mãn $x + y = 6$ và $xy = 5$. Không tính cụ thể giá trị của $x$ và $y$, hãy tính giá trị của biểu thức $T = x^3 + y^3$.\par
\vspace{0.4mm}
\noindent\textit{Bài giải:}\dongke{7}
\\ \hline
\end{tabularx}


\newpage

% #############################################################################
% ========================= TOÁN 8 - TUẦN 03 - BUỔI 02: MÃ ĐỀ 104 ============
% ========================= MẶT TRƯỚC (TRANG 7) ===============================
% #############################################################################

\noindent
\begin{tabularx}{\linewidth}{|X|c|}
\hline
\textbf{PHIẾU KIỂM TRA ĐỊNH KỲ TOÁN 8 - TUẦN 03 - BUỔI 02} & \textbf{Mã đề: 104} \\
\textit{Chủ đề: Ôn tập Chương 1 - Đa thức} & \textit{Thời gian: 45 phút} \\
\hline
\end{tabularx}

\vspace{1mm}

\noindent
\begin{tabularx}{\linewidth}{|X|}
\hline
\rule[-2mm]{0pt}{7.5mm}\textbf{Họ và tên:} \textbf{\large Trần Đình Hải Đăng} \hfill \textbf{Mã HS:} \texttt{HAIDANG2506} \quad \\
\hline
\end{tabularx}

\vspace{1mm}

\noindent
\begin{tabularx}{\linewidth}{|p{3.2cm}|X|}
\hline
\centering\textbf{ĐIỂM SỐ} & \centering\textbf{LỜI NHẬN XÉT CỦA THẦY CÔ} \tabularnewline
\hline
\rule{0pt}{1.3cm} & \tabularnewline
\hline
\end{tabularx}

\vspace{1.2mm}

\noindent\textbf{BẢNG ĐIỀN ĐÁP ÁN TRẮC NGHIỆM (Mỗi câu đúng 0,5 điểm):}
\vspace{0.8mm}

\noindent
\renewcommand{\arraystretch}{1.35}
\begin{tabularx}{\linewidth}{|c|*{5}{>{\centering\arraybackslash}X|}}
\hline
\textbf{Câu} & \textbf{1} & \textbf{2} & \textbf{3} & \textbf{4} & \textbf{5} \\
\hline
\textbf{Đáp án} & & & & & \\
\hline
\end{tabularx}
\renewcommand{\arraystretch}{1.0}

\vspace{1.5mm}

\noindent\textbf{\large PHẦN I. TRẮC NGHIỆM KHÁCH QUAN (2,5 điểm)}
\vspace{1mm}

\cauTN{1}{Bậc của đa thức $Q = 5x^4 y^3 - 2x^3 y^4 + 4x^6 y - 5x^4 y^3 + 7$ là:}{$6$}{$5$}{$7$}{$4$}

\cauTN{2}{Đơn thức nào sau đây đồng dạng với đơn thức $-7x^4 y^2$?}{$3x^4 y^2$}{$-7x^2 y^4$}{$4x^4 y^3$}{$2x^3 y^2$}

\cauTN{3}{Kết quả của phép nhân đa thức $(3x + 2)(x - 2)$ là:}{$3x^2 + 4x - 4$}{$3x^2 - 4x - 4$}{$3x^2 - 6$}{$3x^2 - 8x - 4$}

\cauTN{4}{Kết quả của phép chia $(24x^5 y^3 - 18x^3 y^2) : 6x^3 y^2$ là:}{$4x^2 y - 3$}{$4x^2 y - 3xy$}{$4x^2 y^2 - 3$}{$4x^2 - 3y$}

\cauTN{5}{Giá trị của biểu thức $N = x^2 + 8xy + 16y^2$ tại $x = 12$ và $y = 2$ là:}{$256$}{$400$}{$324$}{$100$}

\vspace{1mm}

\noindent\textbf{\large PHẦN II. TỰ LUẬN (7,5 điểm)}

\noindent\textbf{Câu 1 (2,5 điểm). Thực hiện phép tính với đa thức:}
\vspace{0.8mm}

\noindent
\begin{tabularx}{\linewidth}{|X|X|}
\hline
\multicolumn{2}{|p{0.97\linewidth}|}{Cho hai đa thức: $A = 5x^2 y - 4xy^2 + 3xy - 1$ \ và \ $B = 3x^2 y + 4xy^2 - 2xy + 6$.} \\
\hline
\oBaiTap[7]{a)}{Tính $C = A + B$:} & 
\oBaiTap[7]{b)}{Tính $D = A - B$:} \\
\hline
\oBaiTap[5]{c)}{Tìm bậc của đa thức $C$:} & 
\oBaiTap[5]{d)}{Tính giá trị của đa thức $C$ tại $x = 1; y = -1$:} \\
\hline
\end{tabularx}


\newpage

% #############################################################################
% ========================= TOÁN 8 - TUẦN 03 - BUỔI 02: MÃ ĐỀ 104 ============
% ========================= MẶT SAU (TRANG 8) =================================
% #############################################################################

\noindent
\begin{tabularx}{\linewidth}{|X|c|}
\hline
\textbf{PHẦN II. TỰ LUẬN (Tiếp theo) - TOÁN 8 TUẦN 03 BUỔI 02} & \textbf{Mã đề: 104} \\
\hline
\end{tabularx}

\vspace{1mm}

\noindent
\begin{tabularx}{\linewidth}{|X|}
\hline
\rule[-2mm]{0pt}{7mm}\textbf{Họ và tên:} \textbf{\large Trần Đình Hải Đăng} \hfill \textbf{Mã HS:} \texttt{HAIDANG2506} \quad \\
\hline
\end{tabularx}

\vspace{1.5mm}

\noindent\textbf{\large Câu 2 (2,5 điểm). Rút gọn và tìm $x$:}
\vspace{0.8mm}

\noindent
\begin{tabularx}{\linewidth}{|X|X|}
\hline
\oBaiTap[12]{a)}{Rút gọn và chứng minh giá trị biểu thức không phụ thuộc vào giá trị của biến $x$:\par
$N = (4x - 1)(x + 2) - 4x(x + 1) - 3x + 9$.} & 
\oBaiTap[12]{b)}{Tìm số thực $x$, biết:\par
$(2x - 1)(4x + 3) - 8x(x - 1) = 29$.} \\
\hline
\end{tabularx}

\vspace{1.8mm}

\noindent\textbf{\large Câu 3 (2,5 điểm). Bài toán thực tế và nâng cao:}
\vspace{0.8mm}

\noindent
\begin{tabularx}{\linewidth}{|X|}
\hline
\vspace{0.5mm}
\textbf{a) (1,5 điểm)} Một khu đất hình chữ nhật có chiều dài là $(4x + 5)\text{ m}$ và chiều rộng là $(2x + 4)\text{ m}$. Người ta làm một lối đi xung quanh khu đất rộng $1\text{ m}$, phần đất còn lại bên trong dùng để trồng rau sạch.\par
1) Viết đa thức theo biến $x$ biểu thị diện tích khu đất ban đầu.\par
2) Viết đa thức theo biến $x$ biểu thị diện tích phần đất bên trong dùng để trồng rau.\par
\vspace{0.4mm}
\noindent\textit{Bài giải:}\dongke{9}
\vspace{1mm}
\textbf{b) (1,0 điểm)} Cho hai số $x, y$ thỏa mãn $x + y = 6$ và $xy = 8$. Không tính cụ thể giá trị của $x$ và $y$, hãy tính giá trị của biểu thức $M = x^3 + y^3$.\par
\vspace{0.4mm}
\noindent\textit{Bài giải:}\dongke{7}
\\ \hline
\end{tabularx}



\newpage


% #############################################################################
% ========================= TOÁN 8 - TUẦN 03 - BUỔI 02: MÃ ĐỀ 101 ============
% ========================= MẶT TRƯỚC (TRANG 1) ===============================
% #############################################################################

\noindent
\begin{tabularx}{\linewidth}{|X|c|}
\hline
\textbf{PHIẾU KIỂM TRA ĐỊNH KỲ TOÁN 8 - TUẦN 03 - BUỔI 02} & \textbf{Mã đề: 101} \\
\textit{Chủ đề: Ôn tập Chương 1 - Đa thức} & \textit{Thời gian: 45 phút} \\
\hline
\end{tabularx}

\vspace{1mm}

\noindent
\begin{tabularx}{\linewidth}{|X|}
\hline
\rule[-2mm]{0pt}{7.5mm}\textbf{Họ và tên:} \textbf{\large Võ Minh Đăng} \hfill \textbf{Mã HS:} \texttt{MINHDANG0512} \quad \\
\hline
\end{tabularx}

\vspace{1mm}

\noindent
\begin{tabularx}{\linewidth}{|p{3.2cm}|X|}
\hline
\centering\textbf{ĐIỂM SỐ} & \centering\textbf{LỜI NHẬN XÉT CỦA THẦY CÔ} \tabularnewline
\hline
\rule{0pt}{1.3cm} & \tabularnewline
\hline
\end{tabularx}

\vspace{1.2mm}

\noindent\textbf{BẢNG ĐIỀN ĐÁP ÁN TRẮC NGHIỆM (Mỗi câu đúng 0,5 điểm):}
\vspace{0.8mm}

\noindent
\renewcommand{\arraystretch}{1.35}
\begin{tabularx}{\linewidth}{|c|*{5}{>{\centering\arraybackslash}X|}}
\hline
\textbf{Câu} & \textbf{1} & \textbf{2} & \textbf{3} & \textbf{4} & \textbf{5} \\
\hline
\textbf{Đáp án} & & & & & \\
\hline
\end{tabularx}
\renewcommand{\arraystretch}{1.0}

\vspace{1.5mm}

\noindent\textbf{\large PHẦN I. TRẮC NGHIỆM KHÁCH QUAN (2,5 điểm)}
\vspace{1mm}

\cauTN{1}{Bậc của đa thức $M = 5x^3 y^2 - 2x^2 y^3 + 3x^4 y - 5x^3 y^2 + 1$ là:}{$4$}{$5$}{$6$}{$3$}

\cauTN{2}{Đơn thức nào sau đây đồng dạng với đơn thức $-4x^3 y^2$?}{$2x^2 y^3$}{$-4x^3 y$}{$7x^3 y^2$}{$5x^2 y^2$}

\cauTN{3}{Kết quả của phép nhân đa thức $(2x - 3)(x + 2)$ là:}{$2x^2 + x - 6$}{$2x^2 - x - 6$}{$2x^2 + 7x - 6$}{$2x^2 - 6$}

\cauTN{4}{Kết quả của phép chia $(18x^4 y^3 - 12x^3 y^2) : 6x^2 y^2$ là:}{$3x^2 y - 2xy$}{$3x^2 y - 2$}{$3x^2 y^2 - 2x$}{$3x^2 y - 2x$}

\cauTN{5}{Giá trị của biểu thức $P = x^2 - 6xy + 9y^2$ tại $x = 16$ và $y = 2$ là:}{$100$}{$144$}{$64$}{$400$}

\vspace{1mm}

\noindent\textbf{\large PHẦN II. TỰ LUẬN (7,5 điểm)}

\noindent\textbf{Câu 1 (2,5 điểm). Thực hiện phép tính với đa thức:}
\vspace{0.8mm}

\noindent
\begin{tabularx}{\linewidth}{|X|X|}
\hline
\multicolumn{2}{|p{0.97\linewidth}|}{Cho hai đa thức: $A = 3x^2 y - 2xy^2 + 5xy - 1$ \ và \ $B = 2x^2 y + 2xy^2 - 3xy + 4$.} \\
\hline
\oBaiTap[7]{a)}{Tính $C = A + B$:} & 
\oBaiTap[7]{b)}{Tính $D = A - B$:} \\
\hline
\oBaiTap[5]{c)}{Tìm bậc của đa thức $C$:} & 
\oBaiTap[5]{d)}{Tính giá trị của đa thức $C$ tại $x = 1; y = -1$:} \\
\hline
\end{tabularx}


\newpage

% #############################################################################
% ========================= TOÁN 8 - TUẦN 03 - BUỔI 02: MÃ ĐỀ 101 ============
% ========================= MẶT SAU (TRANG 2) =================================
% #############################################################################

\noindent
\begin{tabularx}{\linewidth}{|X|c|}
\hline
\textbf{PHẦN II. TỰ LUẬN (Tiếp theo) - TOÁN 8 TUẦN 03 BUỔI 02} & \textbf{Mã đề: 101} \\
\hline
\end{tabularx}

\vspace{1mm}

\noindent
\begin{tabularx}{\linewidth}{|X|}
\hline
\rule[-2mm]{0pt}{7mm}\textbf{Họ và tên:} \textbf{\large Võ Minh Đăng} \hfill \textbf{Mã HS:} \texttt{MINHDANG0512} \quad \\
\hline
\end{tabularx}

\vspace{1.5mm}

\noindent\textbf{\large Câu 2 (2,5 điểm). Rút gọn và tìm $x$:}
\vspace{0.8mm}

\noindent
\begin{tabularx}{\linewidth}{|X|X|}
\hline
\oBaiTap[12]{a)}{Rút gọn và chứng minh giá trị biểu thức không phụ thuộc vào giá trị của biến $x$:\par
$M = (2x - 1)(x + 3) - 2x(x + 2) - x + 7$.} & 
\oBaiTap[12]{b)}{Tìm số thực $x$, biết:\par
$(3x - 1)(2x + 1) - 6x(x - 2) = 25$.} \\
\hline
\end{tabularx}

\vspace{1.8mm}

\noindent\textbf{\large Câu 3 (2,5 điểm). Bài toán thực tế và nâng cao:}
\vspace{0.8mm}

\noindent
\begin{tabularx}{\linewidth}{|X|}
\hline
\vspace{0.5mm}
\textbf{a) (1,5 điểm)} Một mảnh vườn hình chữ nhật có chiều dài là $(3x + 4)\text{ m}$ và chiều rộng là $(2x + 1)\text{ m}$. Người ta làm một lối đi xung quanh vườn rộng $1\text{ m}$, phần đất còn lại bên trong dùng để trồng hoa.\par
1) Viết đa thức theo biến $x$ biểu thị diện tích mảnh vườn ban đầu.\par
2) Viết đa thức theo biến $x$ biểu thị diện tích phần đất bên trong dùng để trồng hoa.\par
\vspace{0.4mm}
\noindent\textit{Bài giải:}\dongke{9}
\vspace{1mm}
\textbf{b) (1,0 điểm)} Cho hai số $x, y$ thỏa mãn $x + y = 4$ và $xy = 3$. Không tính cụ thể giá trị của $x$ và $y$, hãy tính giá trị của biểu thức $T = x^3 + y^3$.\par
\vspace{0.4mm}
\noindent\textit{Bài giải:}\dongke{7}
\\ \hline
\end{tabularx}


\end{document}
